%% main_ist.tex --- Information and Software Technology (Elsevier) submission
%% Manuscript source for Information and Software Technology.
%% Class: elsarticle, single-column preprint layout ("Your Paper Your Way").
%% Reference style: elsarticle-num (Elsevier numbered).

\documentclass[preprint,12pt]{elsarticle}

\usepackage[T1]{fontenc}
\usepackage{booktabs}
\usepackage{graphicx}
\usepackage{amsmath}
\usepackage{amssymb}
\usepackage{xcolor}
\usepackage{url}
\usepackage{float}
\usepackage{placeins}
\usepackage{caption}
\usepackage{longtable}
\usepackage{tabularx}
\usepackage{array}
\usepackage{calc}
\usepackage{tikz}
\usetikzlibrary{positioning,arrows.meta}
\usepackage{enumitem}
\usepackage{lineno}
\modulolinenumbers[5]

\providecommand{\tightlist}{\setlength{\itemsep}{0pt}\setlength{\parskip}{0pt}}
\setlist{itemsep=1pt, topsep=2pt, parsep=0pt, leftmargin=1.2em}
\setlength{\tabcolsep}{3pt}
\raggedbottom
\emergencystretch=3em
\hyphenation{re-pro-duc-tion spec-i-fi-ca-tions knowledge-graph-grounded
 taxonomy-grounded An-ten-na-Pod New-Pipe Thun-der-bird
 con-fi-dent cleanlab hi-er-ar-chi-cal cov-er-age IssueSpec}

\journal{Information and Software Technology}

\begin{document}

\begin{frontmatter}

\title{Typed Issue Specifications from App Reviews: What Scales, and a Downstream Benefit That Does Not Replicate}

%% ---------------------------------------------------------------
%% AUTHOR BLOCK. Replace the placeholders before submission.
%% IST is single-blind, so real author names are expected at submission.
%% ---------------------------------------------------------------
\author[inst1]{Fabiha Jalal\corref{cor1}}
\ead{fabihajalal@iut-dhaka.edu}
\author[inst1]{Sadia Tasnim Dhruba}
\ead{sadiatasnim@iut-dhaka.edu}
\author[inst1]{Tazrian Zaman Tushti}
\ead{tazrian@iut-dhaka.edu}
\author[inst2]{Ajwad Abrar}
\ead{ajwadabrar@iut-dhaka.edu}
\author[inst2]{Md Kamrul Hasan}
\ead{hasank@iut-dhaka.edu}
\author[inst2]{Hasan Mahmud}
\ead{hasan@iut-dhaka.edu}

\cortext[cor1]{Corresponding author.}

\affiliation[inst1]{organization={Department of Computer Science, Islamic University of Technology (IUT)},
            city={Gazipur},
            country={Bangladesh}}
\affiliation[inst2]{organization={Systems and Software Lab (SSL), Department of Computer Science and Engineering, Islamic University of Technology (IUT)},
            city={Gazipur},
            country={Bangladesh}}

%% ---------------------------------------------------------------
%% STRUCTURED ABSTRACT. IST requires the five headings below; papers
%% without them are not handled by the editorial office.
%% ---------------------------------------------------------------
\begin{abstract}
\textbf{Context.} Developers receive millions of noisy app reviews, but defect trackers such as
JIRA and GitHub Issues cannot consume free text. They require structured fields such as issue
type, steps to reproduce, severity, and affected component. Review-mining research stops short of
that boundary, predicting labels or writing replies without producing the record a tracker can
ingest.

\textbf{Objective.} Can a typed intermediate representation (IR) of a review be produced at
corpus scale from noisy labels, and does supplying it to a response generator improve the reply a
user receives? We answer the first yes and the second no, and the second answer corrects a claim
we made ourselves.

\textbf{Method.} IssueSpec routes each review by issue class into one of five standards-body
templates, repairs the noisy labels the routing depends on with a verified-anchor procedure, and
feeds the result to a retrieval-augmented generator. We study 215,583 reviews from 58 Android
apps with a 5,230-review verified anchor, a 490-review three-rater gold standard, six ablations,
and four language models. The downstream evaluation was run twice: first with conditions written
by different template composers, then with the generator held fixed so that only the IR varies,
all 200 blinded rows scored by three raters.

\textbf{Results.} Production works. The verified-anchor procedure lifts Cohen's $\kappa$ against
the human gold from 0.163 to 0.592, and to 0.616 on reviews the classifier never saw. Type
routing reaches 0.96 template-fill against 0.69 without it. The downstream benefit does not. The
first round measured $+2.35$ Likert points, but its arms differed in authorial register as much
as in information, one emitting lemmatised text at 78 words and the other polished prose at 123;
the corrected comparison gives $+0.03$ ($p = 0.38$), with raters disagreeing in direction. Three
further components fail their own tests: the knowledge-graph layer loses to count-matched flat
clustering, our extractive-coverage faithfulness scorer shows no rank correlation with human
judgement, and the constrained alignment layer never demonstrates the trade-off it was built to
expose.

\textbf{Conclusions.} Typed issue specifications are producible from noisy reviews at corpus
scale, and that is the contribution we defend. Their claimed downstream payoff is not. We release
the arithmetic check that distinguishes our two rating rounds, because a structured intermediate
should be justified on what it produces rather than on an effect a fair comparison removes.
\end{abstract}

%% ---------------------------------------------------------------
%% HIGHLIGHTS. Elsevier allows 3-5 bullets, each at most 85 characters.
%% Also supply these in a separate "highlights.tex"/"Highlights" file at
%% submission; the block below is for the reviewer-facing PDF.
%% ---------------------------------------------------------------
\begin{highlights}
\item Typed intermediate representation turns noisy app reviews into tracker-ready issues.
\item A verified-anchor procedure lifts Cohen's kappa from 0.163 to 0.592 at corpus scale.
\item Template-field fill rate reaches 0.96 against 0.69 without type routing.
\item A clean same-model re-run finds no downstream response-quality gain from the IR.
\item We release the check that distinguishes independent ratings from derived ones.
\end{highlights}

\begin{keyword}
App review mining \sep Issue report generation \sep Intermediate representation \sep
Knowledge graph \sep Retrieval-augmented generation \sep Reinforcement learning from human
feedback \sep Empirical software engineering
\end{keyword}

\end{frontmatter}

\linenumbers

\section{Introduction}\label{introduction}

App stores are where users talk back to developers. A single popular app pulls in thousands of reviews a week, and the global stream runs into millions a day~\cite{maalej2015bug, gao2019automating, dkabrowski2022analysing}. These reviews carry the signals developers need to act on: bug reports (a crash on Samsung S8 after an update), feature requests (dark mode), performance complaints (battery drain), and regression notes (voice search no longer works). The information is real, but it shows up as messy free-form text, usually without the details required to reproduce a bug or scope a new feature.

To turn any of this into work, developers rely on defect trackers such as JIRA or GitHub Issues. Those tools, though, can't work with raw text; they need structured fields such as issue type, steps to reproduce, severity, and affected component~\cite{chaparro2017detecting,
zimmermann2010makes}. Filling those fields by hand does not scale with the volume modern apps face. Years of work on automated review pipelines has helped with the surrounding tasks like classification, summarization, and response generation, but no system actually produces the structured report a developer can paste directly into a tracker. D{\k a}browski et al.'s nineteen-year survey of the field~\cite{dkabrowski2022analysing} confirms the same boundary.

Prior work has tackled several pieces of this problem separately.
  Classifiers~\cite{maalej2015bug, chen2014ar, villarroel2016release, di2017surf} take a review
  and predict an issue label such as bug or feature request. Response
  generators~\cite{gao2019automating, gao2021automating, lewis2020retrieval} produce a polite reply for the user.
   Agentic software-engineering tools such as SWE-Agent~\cite{yang2024swe} and
  RepairAgent~\cite{bouzenia2025repairagent} take a clean GitHub issue and generate a code
  patch. Knowledge-graph methods~\cite{de2025reviewgraph, yu2024knowledge, keertipati2016approaches} cluster
  reviews into themes or surface frequent complaints. Each of these has matured over the past
  decade and works well within its own scope.



But none of these reaches the developer-ready issue report a defect tracker needs. The
  Maalej-Nabil classifier taxonomy~\cite{maalej2015bug} and its successors~\cite{chen2014ar,
  villarroel2016release, di2017surf, kaur2023bert, shah2020mining} have settled at an F1
   ceiling around 0.75--0.85, a label-quality limit (not a model-capacity one) per D{\k
  a}browski et al.'s~\cite{dkabrowski2022analysing} nineteen-year survey. The label is only the
  first piece of what a tracker entry needs; the remaining fields (steps to reproduce, severity,
   affected component) have to come from somewhere, and classification-only pipelines leave that
   downstream step to manual triage. Response generators like RRGen~\cite{gao2019automating} and
  CoRe~\cite{gao2021automating} address the customer-communication side of the workflow, but the
  actual bug behind the review is never written out in the structured form a tracker needs, so
  even a perfect reply system does not give developers something they can paste into JIRA.
  Agentic SE tools assume a well-formed bug report or GitHub issue as the agent's starting
  state, built into the SWE-bench protocol~\cite{jimenez2023swe} where each instance is a
  (GitHub issue, code repository) pair; constructing such an issue from raw review text is
  explicitly outside the protocol, leaving the upstream gap unaddressed.

Knowledge-graph methods stop at intermediate outputs rather than typed fields:
  ReviewGraph~\cite{de2025reviewgraph}, KGCPN~\cite{yu2024knowledge}, and Keertipati et
  al.~\cite{keertipati2016approaches} each output rating predictions, summaries, or feature rankings, none
   of which a defect tracker can ingest directly. RLHF has matured into the standard approach
  for aligning LLM outputs with human preferences (DPO~\cite{rafailov2023direct},
  KTO~\cite{ethayarajh2024kto}, PPO~\cite{schulman2017proximal}), and Safe-RLHF~\cite{dai2024safe} extends
  the framework to multi-objective settings under a Constrained MDP for general LLM safety
  (harmfulness, dishonesty, bias). The developer-relations setting poses a different question
  entirely: how to align responses under a domain-specific constraint family (truthful promises,
   no information leakage, professional tone, legal-safe wording) rooted in business and legal
  exposure. To our knowledge, no published RLHF formulation has been adapted to this constraint
  set.

This work pursues three aims. First, we develop a framework that turns noisy app reviews into structured, taxonomy-grounded issue specifications a defect tracker can consume directly. Second, we build a response generator that uses these specifications as input, producing replies that reference the actual underlying defect rather than generic templates. Third, we close the loop with a dual-objective alignment trainer that jointly optimizes for response helpfulness and the domain-specific compliance constraints of the developer-relations setting.

These aims map to three research questions.

\textbf{RQ1 (Translation Quality).} Can a typed IR, generated via a three-layer framework (knowledge-graph construction, hierarchical clustering,
standardized schema mapping), produce issue specifications that match or exceed the substantive completeness of human-written GitHub issues?

\textbf{RQ2 (Downstream Load-Bearing).} Does conditioning a response generator on the validated IssueSpec measurably improve human-rated response quality over RAG alone?

\textbf{RQ3 (Dual-Objective Alignment).} Can a CMDP-grounded RLHF loop progressing through KTO, DPO, and Constrained PPO jointly optimize response helpfulness and policy compliance for the
developer-relations setting? We answer this one in the negative for the scale we can train at, and report what that rules out.

To address these challenges, we propose \textbf{IssueSpec}, a framework for structured review-to-issue translation. It consists of two parts: a typed Intermediate Representation (IR) that converts each noisy review into one of five standard report types, and a CMDP-grounded response generator that consumes the IR under quality and compliance constraints. Figure \ref{fig:architecture} shows the full five-stage pipeline.

\textbf{What is new, and what is assembled.} Several components of this pipeline are existing
technique used as intended: RoBERTa for multi-label classification, confident
learning~\cite{northcutt2021confident} for label noise, Sentence-BERT with UMAP and HDBSCAN for
clustering, PageRank for centrality, retrieval-augmented generation for response drafting, and
the CMDP formulation of Altman~\cite{altman2021constrained}. We do not claim any of these. Three things are
new. \emph{First}, the typed intermediate representation itself: a review is routed by predicted
issue class into one of five standards-body templates, so the output is a typed record with
per-type obligations rather than a label or a summary. No prior review-mining system emits a
structure a defect tracker can ingest, which Table~\ref{tab:cross-category} makes concrete
against representative systems. \emph{Second}, the verified-anchor procedure, which is what makes
that representation producible from noisy LLM labels at corpus scale rather than on a hand-curated
sample, and which turns out to work because the noise is structurally biased rather than random.
\emph{Third}, the isolation result: we show the typed intermediate, not retrieval and not
iteration, is what carries downstream response quality, and we show it by removing each in turn.
The assembly is not the contribution; the typed object in the middle of it is.

\textbf{The central claim, stated once.} A pipeline with five stages invites the question of
which one the paper is actually about. The answer is Stage~3 and the object it produces: a
\emph{typed, type-routed intermediate representation} of an app review, together with the
verified-anchor procedure that makes it producible at corpus scale from noisy labels. Everything
else in the pipeline exists to feed that object or to consume it. Two experiments carry the
claim. Removing type routing drops strict template-fill from 0.96 to 0.69
(Section~\ref{exp1-results}). We had also claimed a downstream payoff, and Section~\ref{exp2-results}
reports why we withdrew it: under a comparison that holds the generator fixed, supplying the IR
changes response quality by $+0.03$ Likert points, which is nothing. Section~\ref{load-bearing} states which components did
not earn their place on our own evidence, including two we withdraw. A reader who wants the
claim and not the machinery can read Sections~\ref{exp1-results}, \ref{exp2-results},
and~\ref{load-bearing} alone.

The contributions are stated in two groups, because the honest account of this work has a
positive half and a negative half, and separating them is more useful than blending them.

\emph{What we establish.}
\begin{itemize}\setlength{\itemsep}{2pt}
  \item \textbf{A typed intermediate representation for review-to-issue translation.} Each review
  is routed by predicted issue class into one of five standards-body templates, so the output is
  a record a defect tracker can ingest rather than a label or a summary. No prior review-mining
  system emits one (Table~\ref{tab:cross-category}).
  \item \textbf{A verified-anchor procedure that makes the IR producible from noisy labels at
  corpus scale.} Confident learning against a 5{,}230-review human anchor lifts Cohen's $\kappa$
  against a three-rater human gold from 0.163 to 0.592, and to 0.616 on the reviews the
  classifier never saw. It works because the LLM annotation noise is concentrated on minority
  classes rather than spread uniformly, which is a finding about LLM-as-annotator failure that
  transfers beyond this task.
\end{itemize}

\emph{What we withdraw or report as negative.}
\begin{itemize}\setlength{\itemsep}{2pt}
  \item \textbf{The downstream payoff of the IR.} We previously measured $+2.35$ Likert points of
  response-quality gain from supplying the IR. That comparison was confounded: its two arms were
  written by different template composers, not by one model. Holding the generator fixed gives
  $+0.03$ ($p = 0.38$). We report the correction, the two rating rounds, and the check that
  distinguishes them (Section~\ref{exp2-results}).
  \item \textbf{Three components that do not earn their place.} The knowledge-graph layer loses
  to count-matched flat clustering on every intrinsic metric; our extractive-coverage
  faithfulness scorer shows no rank correlation with human judgement and orders generators
  backwards; the CMDP alignment layer engages once the policy can violate its constraint but
  never demonstrates a quality-versus-compliance trade-off.
  \item \textbf{A reusable check on rating provenance.} Ratings derived from another rater's
  vector by small perturbation leave an arithmetic signature: deviation sets that do not overlap
  and a single deviation magnitude. Our discarded round shows it and our reported round does not.
  We release both so the check can be applied to other annotation artifacts, including ours.
\end{itemize}

\begin{table}[!htbp]\centering\footnotesize
\caption{Contribution-to-evidence-to-artifact map. Each contribution states a measurable headline result, the table or section that anchors it, and the named artifact released with this paper.}\label{tab:c2evidence}
\begin{tabular}{@{}p{1.5cm}p{3.0cm}p{2.9cm}@{}}\toprule
contribution & headline evidence & released artifact \\\midrule
IssueSpec & 0.96 template-fill vs 0.53 GitHub (Table~\ref{tab:stage3}); 5-dim rubric 3.94/5 on all 100 specs (Table~\ref{tab:rubric}) & type-routed schema (Table~\ref{tab:templates}) \\
\textsc{SpecCov} (negative) & no rank correlation with human faithfulness ($\rho=0.15$, $p=0.26$); orders generators backwards (\S\ref{speccov-validation}) & the scorer plus 60 specs rated by three humans \\
IR downstream payoff (withdrawn) & the $+2.35$ gain came from a confounded design; a same-model comparison gives $+0.03$ ($p=0.38$) with raters disagreeing in direction (\S\ref{exp2-results}) & both rating rounds, the generated replies, and the independence check \\
verified anchor & $\kappa$ 0.16$\to$0.59 (Table~\ref{tab:kappa}); per-class F1 recovery
  (Table~\ref{tab:perclass-f1}); $\kappa = 0.616$ on the 307 strictly held-out gold reviews & 5{,}230-review
  verified anchor + cleanlab correction procedure + V5 classifier \\
KG prioritization & $5.4\times$ lower Davies-Bouldin (Table~\ref{tab:cluster-quality}); top-5 PageRank aspects cover 74\% of aspect mentions and 55\% of reviews; count-controlled A1b (\S\ref{ablations}) & three-layer KG construction + count-controlled A1b output \\
CMDP RLHF testbed (partial) & constraint never binds at distilGPT2 scale; with a policy able to
  violate it binds on 13/90 steps and $\lambda$ rises monotonically to 0.65, with quality and
  compliance measurably in tension, but the loop back to compliance is not closed
  (\S\ref{exp3-results}) & CMDP formulation, KTO/DPO/constrained-proxy/CPPO trainers, capability
  probe, binding-constraint re-run, threshold sweep, benchmark \\
\bottomrule\end{tabular}\end{table}

 \begin{figure}[t]\centering
  \includegraphics[width=\textwidth]{fig_architecture.pdf}
  \caption{The IssueSpec five-stage pipeline. Reviews flow left to right through intake and
  classification (Stage~1), three-layer knowledge-graph clustering (Stage~2), LLM-agent IR
  translation (Stage~3), spec-aware RAG (Stage~4), and CMDP-grounded RLHF (Stage~5), each
  emitting a typed intermediate. Yellow checkpoints mark human-in-the-loop review under
  selective prediction~\cite{geifman2017selective, bansal2019beyond}; the red loop propagates Stage~5
  corrections upstream.}\label{fig:architecture}
  \end{figure}
%% =====================================================================
\section{Related Work}\label{related-work}

We organize related work by the five gaps identified in Section~\ref{introduction}, naming what each field does, what assumption breaks for review-to-issue translation, and how IssueSpec differs.

\subsection{App Review Classification and Mining}
Maalej and Nabil~\cite{maalej2015bug} established the canonical seven-class taxonomy for app reviews. Subsequent systems, AR-Miner~\cite{chen2014ar}, CLAP~\cite{villarroel2016release}, SURF~\cite{di2017surf}, BERT/RCNN hybrids~\cite{kaur2023bert, shah2020mining}, and per-aspect sentiment analysis~\cite{guzman2014users}, refined classification accuracy and added aspect-level granularity. Dąbrowski et al.~\cite{dkabrowski2022analysing} surveyed nineteen years of such work and concluded that label quality, not model capacity, is the F1 ceiling at 0.75--0.85, an observation we corroborate in Section~\ref{exp1-results}. \emph{What breaks for review-to-issue translation:} none of these systems emits a typed developer-routable artifact downstream of the classification label. IssueSpec differs by routing the predicted class to one of five typed templates rather than terminating at the label.

\subsection{Response Generation and Retrieval-Augmented Generation}
RRGen~\cite{gao2019automating} and CoRe~\cite{gao2021automating} pioneered end-to-end neural response generation on (review, response) pairs without an explicit structured intermediate. Vanilla RAG~\cite{lewis2020retrieval} retrieves relevant past responses in a single pass. \emph{Agentic RAG} variants, Self-RAG~\cite{asai2024self} and DEMONSTRATE-SEARCH-PREDICT (DSP)~\cite{khattab2022demonstrate}, introduce iterative critique-revise loops that, in principle, could close the structured-intermediate gap. \emph{What breaks for review-to-issue translation:} none of these systems names the underlying defect, and our $n=10$ feasibility study (Section~\ref{exp2-results}) shows that agentic iteration without a typed intermediate adds modest gains (0.58$\to$0.70 rule-based score) compared with the IssueSpec, whose downstream gain we measured at $+0.03$ Likert points once the generator was held fixed (Section~\ref{exp2-results}). IssueSpec differs by injecting the typed intermediate as a first-class RAG source alongside retrieved past responses, and the contribution we defend is the specification itself rather than its effect on the reply.

\subsection{Agentic Software Engineering}
SWE-Agent~\cite{yang2024swe}, RepairAgent~\cite{bouzenia2025repairagent}, and HyperAgent~\cite{phan2024hyperagent} consume well-structured GitHub issues to produce code patches through multi-step planning and tool use. \emph{What breaks for review-to-issue translation:} the issue specification is the input contract for all three systems; how a defect-tracker entry is constructed from noisy user-facing text is out of scope. IssueSpec is precisely the artifact these systems are designed to consume.

\subsection{Knowledge Graphs for App Reviews and Software Engineering}
Knowledge graphs have been used in software engineering for bug localization, report deduplication, and code understanding, but the issue specification is invariably treated as input rather than output. For app reviews specifically, ReviewGraph~\cite{de2025reviewgraph} learns KG embeddings for rating prediction, KGCPN~\cite{yu2024knowledge} produces KG-based extractive summaries, and Keertipati et al.~\cite{keertipati2016approaches} apply graph centrality for coarse feature prioritization. \emph{What breaks for review-to-issue translation:} none of these systems maps KG nodes to a typed, developer-actionable issue schema, and no prior KG work has been applied to the specific task of producing GitHub-issue-quality artifacts from review text. IssueSpec's three-layer KG pipeline, structured-object layer, aspect-grounded hierarchical clustering layer, and PageRank-prioritized issue-mapping layer, is, to our knowledge, the first KG architecture that emits a typed IR for review-to-issue translation.

\subsection{RLHF and Constrained Optimization for Response Generation}
DPO~\cite{rafailov2023direct}, KTO~\cite{ethayarajh2024kto}, and PPO~\cite{schulman2017proximal} compress quality and compliance signals into a single scalar reward. Altman's CMDP framework~\cite{altman2021constrained} separates these objectives, and Safe-RLHF~\cite{dai2024safe} applies CMDP to general LLM safety with a domain-agnostic compliance set covering harmful, dishonest, and biased outputs. \emph{What breaks for review-to-issue translation:} no prior work applies CMDP to dev-rel response generation, where the compliance set is domain-specific (promise safety, information leakage, tone guidelines, and legal safety). Our Stage~5 formulation differs by instantiating these four domain-specific constraints inside the CMDP Lagrangian and releasing a proof-of-concept training loop.

\subsection{Cross-Category Positioning}
Table~\ref{tab:cross-category} positions IssueSpec against representative prior systems across the five pipeline stages. Existing systems occupy at most two of the five stages; composing all five with a standardized issue schema, three theoretically-grounded HITL checkpoints, KG-centrality prioritization, and a CMDP-grounded RLHF formulation is the structural contribution of this work.

\begin{table}[t]
  \centering
  \footnotesize
  \setlength{\tabcolsep}{4pt}
  \caption{Cross-category positioning. Existing systems compose at most two of the five pipeline
   stages. IssueSpec is the first to compose all five with a standardized issue schema, three
  HITL checkpoints grounded in selective prediction theory, KG-centrality prioritization, and a
  CMDP-grounded RLHF formulation. ``--'' denotes a stage absent from that system.}
  \label{tab:cross-category}
  \begin{tabular}{lccccc}
  \toprule
  System & Classify & Cluster & IssueSpec & Response & RLHF \\
  \midrule
  AR-Miner~\cite{chen2014ar} & filter & LDA topic & -- & -- & -- \\
  CLAP~\cite{villarroel2016release} & rating & DBSCAN & -- & -- & -- \\
  SURF~\cite{di2017surf} & intention & LDA & -- & summary & -- \\
  RRGen~\cite{gao2019automating} & -- & -- & -- & seq2seq & -- \\
  CoRe~\cite{gao2021automating} & -- & -- & -- & context-aware retrieval & -- \\
  Maalej~\cite{maalej2015bug} & 7-class & -- & -- & -- & -- \\
  ReviewGraph~\cite{de2025reviewgraph} & rating & KG & -- & -- & -- \\
  SWE-Agent~\cite{yang2024swe} & -- & -- & consumes (not produces) & -- & -- \\
  Safe-RLHF~\cite{dai2024safe} & -- & -- & -- & general & dual-CMDP (general) \\
  \textbf{IssueSpec (ours)} & \textbf{V5 corrected} & \textbf{KG + PageRank} &
  \textbf{type-routed schema} & \textbf{spec-aware RAG} & \textbf{dual-CMDP (dev-rel)} \\
  \bottomrule
  \end{tabular}
  \end{table}

%% =====================================================================
\section{The IssueSpec Pipeline}\label{system-design}

The IssueSpec pipeline (Fig.~\ref{fig:architecture}) is a five-stage information-extraction cascade~\cite{hearst1999untangling, sarawagi2008information}. Stages~1--3 produce the typed IR from raw reviews. Stage~4 uses it to drive spec-aware response generation. Stage~5 aligns the response generator via constrained RLHF. We describe each stage in turn, with the verified-by-citation taxonomy templates in Table~\ref{tab:templates}. Table~\ref{tab:acronyms} collects every abbreviation used in the paper, so a reader can enter at any section. This section describes the framework only. Evaluation design, baselines, and metrics are in Section~\ref{experiments}, exact settings and hyperparameters are in Appendix~\ref{appendix:settings}, and results are in Section~\ref{results}.

Before the stages themselves, Table~\ref{tab:why-stages} answers the question a five-stage
pipeline invites: what forces each stage to exist, and what breaks without it. We answer it with
the ablation that removes each one, so the justification is measured rather than asserted, and
two of the five turn out not to justify themselves.

\begin{table}[t]\centering\footnotesize
\caption{Why each stage exists, and what the ablation says when it is removed. Two stages do not
justify themselves on our own evidence, and the paper says so rather than defending
them.}\label{tab:why-stages}
\begin{tabular}{@{}p{1.5cm}p{3.4cm}p{3.4cm}p{2.6cm}@{}}\toprule
stage & the problem it must solve & why it cannot be skipped & measured verdict \\\midrule
1 Intake & the corpus label is wrong 25\% of the time, concentrated on minority classes & Stage~3 routes on the issue type, so a wrong type produces the wrong template & \textbf{justified}: $\kappa$ $0.163 \to 0.592$ \\
2 KG & 215{,}583 reviews cannot be triaged one at a time; developers need an order & a spec is written per cluster, so reviews must be grouped and ranked & \textbf{partly}: prioritization yes, cluster quality no (A1b, A2) \\
3 IR & trackers need typed fields, not prose & this is where the typed record is produced & \textbf{justified}: fill 0.96 vs 0.69 without routing (A3) \\
4 RAG & a reply needs house voice and precedent, not just the defect & supplies phrasing the spec does not carry & \textbf{weakly}: removing it moves metrics $<1\%$ (A5) \\
5 CMDP & replies must not over-promise or leak internals & compliance is a hard constraint, not a reward term & \textbf{partly}: constraint binds, trade-off unshown \\
\bottomrule\end{tabular}\end{table}

\begin{table}[t]\centering\footnotesize
\caption{Abbreviations used in this paper, with the section that introduces each.}\label{tab:acronyms}
\begin{tabular}{@{}llc@{}}\toprule
term & expansion & first used \\\midrule
IR & Intermediate Representation, the typed issue spec & \S\ref{introduction} \\
IssueSpec & one filled instance of the typed IR & \S\ref{introduction} \\
KG & Knowledge Graph, the three-layer review index & \S\ref{aim1} \\
HITL & Human-In-The-Loop checkpoint (HITL-1 to HITL-3) & \S\ref{aim1} \\
\textsc{SpecCov} & extractive-coverage score for a spec (withdrawn as a faithfulness measure, \S\ref{speccov-validation}) & \S\ref{aim1} \\
RAG & Retrieval-Augmented Generation & \S\ref{aim2} \\
RLHF & Reinforcement Learning from Human Feedback & \S\ref{aim3} \\
CMDP & Constrained Markov Decision Process & \S\ref{aim3} \\
KTO & Kahneman-Tversky Optimization, a preference trainer & \S\ref{aim3} \\
DPO & Direct Preference Optimization & \S\ref{aim3} \\
PPO & Proximal Policy Optimization; CPPO is its constrained form & \S\ref{aim3} \\
SFT & Supervised Fine-Tuning, the RLHF starting policy & \S\ref{aim3} \\
NFR & Non-Functional Requirement (ISO/IEC~25010) & \S\ref{aim1} \\
DB / CH & Davies-Bouldin / Calinski-Harabasz cluster indices & \S\ref{cluster-validation-results} \\
\bottomrule\end{tabular}\end{table}

\subsection{Producing the IR}\label{aim1}

\textbf{Stage 1: Intake, Classification, and Aspect Extraction.}
  Stage~1 produces three outputs per review. The first is a multi-label class from a RoBERTa
  head fine-tuned on Maalej~\cite{maalej2015bug} and Guzman~\cite{guzman2014users}. Labels are
  multi-label because reviews mix concerns, as when a bug report also asks for a feature. The
  second is aspect-level sentiment~\cite{guzman2014users}, which keeps ``great UI but terrible
  battery'' as two separate pairs instead of averaging it into one star rating. The third is
  entity extraction for devices, OS versions, and named features. Stage~3 needs those entities
  as typed fields for the device-OS and bug-report templates.
  Low-confidence reviews route to HITL-1 under selective-prediction
  theory~\cite{geifman2017selective}: when the classifier cannot decide,
  the expert does, and the correction feeds the next training round.

  LLM annotation noise is the main challenge here. On the
  5{,}041-review praise stratum, 25\% of LLM labels are wrong,
  concentrated on minority classes. Discarding them wastes annotation;
  using them as-is poisons training. We apply confident
  learning~\cite{northcutt2021confident} against a verified anchor of
  5{,}230 expert-labeled reviews plus 5{,}008 from
  Maalej~\cite{maalej2015bug}, gated by a confidence threshold on the anchor and a probability
  ceiling on the LLM label (Appendix~\ref{appendix:settings}); this keeps the good labels and
  corrects only the wrong ones. The compatibility class needs a separate fix:
  only 8 of 215{,}583 LLM labels are tagged \texttt{compatibility}, and
  Maalej has none. We seed it with template-generated synthetic
  reviews~\cite{joorabchi2013real} and keyword-mined RRGen reviews, each checked by the lead
  author before training; counts are in Appendix~\ref{appendix:settings}.
  
  \textbf{Stage 2: Three-Layer Knowledge Graph.} Stage~2 organizes
  reviews into an aspect-grounded knowledge graph: build, cluster,
  prioritize. The first layer builds the graph: each review is a node
  linked to its aspects and entities by sentiment and temporal edges. A
  graph beats a flat list because reviews share aspects across users,
  and that co-occurrence is what later layers cluster and rank on.
  Aspects come from dependency-parse noun phrases following Hu and
  Liu~\cite{hu2004mining}, pruned by frequency and centrality to drop
  one-off terms.

  The second layer clusters reviews first by aspect, then sub-clusters
  within each aspect using sentence
  embeddings~\cite{reimers2019sentence}, UMAP~\cite{mcinnes2018umap},
  and HDBSCAN~\cite{mcinnes2017hdbscan}. Aspect-first keeps each
  sub-cluster on one theme: a flat clusterer mixes ``battery on Samsung''
  with ``battery in dark mode'' because both share lexical content. The
  result is 605 aspect-grounded sub-clusters with mean size 3.0, against
  194 from a flat baseline. The third layer prioritizes via
  PageRank~\cite{brin1998anatomy, page1999pagerank} on the bipartite
  review-aspect graph, ranking aspects by cross-review prevalence so
  developers see the most-affected first. Each
  cluster maps to the fixed schema in Table~\ref{tab:templates} with a
  centrality-derived priority score. The resulting aspect ranking is reported in
  Section~\ref{cluster-validation-results}.

  \textbf{Stage 3: Review-to-Issue Translation.} Stage~3 instantiates
  the IssueSpec. An LLM agent reads the standardized cluster schema and
  writes a GitHub-issue-quality specification using the type-routed
  templates in Table~\ref{tab:templates}. We use an LLM rather than
  rule-based extraction because reviews vary too much for fixed rules
  and the agent has to reason across a 5-review cluster. We route by
  issue type because a bug report needs reproduction steps, a feature
  request needs acceptance criteria, and a performance complaint needs
  latency or memory numbers; a shared template would be incomplete or
  padded. The five templates are Zimmermann bug
  reports~\cite{zimmermann2010makes}, Cohn user stories~\cite{cohn2004user}
  with BDD acceptance criteria~\cite{hellesoy2017cucumber}, ISO/IEC~25010
  non-functional categories~\cite{international2011systems}, Nielsen
  heuristics~\cite{nielsen1994usability}, and device-OS
  matrices~\cite{joorabchi2013real}. Each is grounded in its
  standards-body citation because developers already recognize these
  formats.
  
  The output passes through HITL-2: a domain expert scores each spec on
  a five-dimension rubric (completeness, accuracy, actionability,
  specificity, clarity, each 1--5), following the LLM-as-judge
  calibration convention of Gilardi et al.~\cite{gilardi2023chatgpt},
  Pangakis et al.~\cite{pangakis2023automated}, and Zheng et
  al.~\cite{zheng2023judging}. Five dimensions, not one, because a
  single quality number conflates completeness with clarity and the LLM
  cannot tell which one to fix on retry. Any spec failing on any
  dimension goes back to the LLM with dimension-level feedback for a
  second pass.


\textbf{\textsc{SpecCov} faithfulness metric.} We release \textsc{SpecCov}, an extractive-coverage scorer that measures substantive-token overlap (length~$\geq 5$, stop-word filtered) with quoted-string and coverage bonuses. Bands follow news-summarization conventions~\cite{grusky2018newsroom}; the metric targets the extractive-coverage sub-type of faithfulness~\cite{grusky2018newsroom, maynez2020faithfulness}, a known lexical correlate of factual consistency~\cite{kryscinski2020evaluating} (the remaining hallucination sub-types in Ji et al.~\cite{ji2023survey} are out of scope). Scores are reported with the other Stage~3 results in Section~\ref{exp1-results}.
\textbf{What \textsc{SpecCov} is and is not.} It is a lexical-grounding proxy in the sense
of~\cite{grusky2018newsroom}: it measures whether the content of a spec can be traced back to
tokens in its source cluster. It is not a faithfulness classifier, and Section~\ref{speccov-validation}
reports the experiment that establishes this. We describe it here because it is part of the
released pipeline and because the negative result only makes sense against its definition.

\begin{table}[t]
  \centering
  \footnotesize
  \caption{IssueSpec standardized issue schema. Each issue type follows the dominant industry or
   HCI standard; templates are verified against their standards-body citations, producing
  GitHub-issue-quality artifacts a defect tracker can consume directly.}\label{tab:templates}
  \begin{tabular}{@{}p{2.2cm}p{8.5cm}p{6.0cm}@{}}
  \toprule
  Issue type & Template (citation) & Justification \\
  \midrule
  \texttt{bug\_report} & Zimmermann~\cite{zimmermann2010makes}: title, severity,
  steps\_to\_reproduce, expected, actual, component & Standard in GitHub Issues, JIRA, defect
  trackers; Chaparro et al.~\cite{chaparro2017detecting} confirm these are the most-requested
  bug-report fields \\
  \texttt{feature\_request} & Cohn user-story~\cite{cohn2004user} (As-a / I-want / So-that) +
  behavior-driven development (BDD) acceptance criteria~\cite{hellesoy2017cucumber} & Dominant
  industry format (Agile, Scrum, SAFe); BDD standardizes testable outcomes \\
  \texttt{performance} & ISO/IEC~25010~\cite{international2011systems} performance-efficiency
  sub-characteristics: time behavior (speed, responsiveness), resource utilization (memory,
  battery), capacity (scalability) & International standard for software-quality NFR taxonomy;
  the parenthesized terms map ISO sub-characteristics to mobile-app vocabulary \\
  \texttt{usability} & Nielsen-10 heuristics~\cite{nielsen1994usability} (visibility, match-real-world,
  user control, etc.) & Most-cited usability evaluation framework in human-computer interaction
  \\ 
  \texttt{compatibility} & Device-OS matrix (devices $\times$ OS
  versions)~\cite{joorabchi2013real} & Device/OS fragmentation is the canonical mobile-QA
  reporting axis \\   
  \bottomrule
  \end{tabular}
  \end{table}

\subsection{Using the IR in Response Generation}\label{aim2}

Stage~4 is a RAG layer over five fixed sources, 15{,}100 documents in total in the released
  ChromaDB index. Four of them are pre-indexed: app changelogs and release notes (60 docs), FAQ
  and help documentation (40), 10{,}000 past review-response pairs from
  RRGen~\cite{gao2019automating}, and 5{,}000 semantically similar past responses. The fifth is the
  validated IssueSpec from Stage~3, injected per query rather than pre-indexed, because specs
  are generated per cluster. These five were chosen to give the generator the answer (changelogs, FAQ), the voice (past responses), the closest precedent (semantic neighbors), and the structured intent (IssueSpec) without overlap.
  
\subsection{Aligning the IR-Aware Generator}\label{aim3}

Stage~5 frames alignment as a CMDP~\cite{altman2021constrained, dai2024safe}:
  maximize a quality reward subject to a compliance constraint,
  \[
  \max_\theta \mathbb{E}_{\pi_\theta}[R_{\text{quality}}] \quad
  \text{s.t.} \quad \mathbb{E}_{\pi_\theta}[C_{\text{compliance}}] \le
  \tau,
  \] 
  with the Lagrangian relaxation
  \[
  \mathcal{L}(\theta,\lambda) = \mathbb{E}_{\pi_\theta}[R - \lambda(C -
  \tau)]    
  \]
  collapsing gracefully to single-objective optimization whenever the
  constraint is already satisfied. We use a CMDP rather than a scalar
  reward because a dev-rel response can be high-helpfulness and still
  out of compliance, like a polite promise the team cannot deliver or a
  leaked roadmap detail, and a weighted sum hides exactly the trade-off
  the constraint is meant to expose. The compliance scorer tracks four
  domain-specific dimensions absent from prior general-safety
  RLHF~\cite{dai2024safe}: promise safety (no commitments the team cannot
  keep), information leakage, tone guidelines, and legal safety. The
  trainer is Constrained PPO~\cite{schulman2017proximal, dai2024safe} with a
  KL-regularized reference policy that keeps updates from drifting too
  far from the supervised baseline.
  KTO~\cite{ethayarajh2024kto} and DPO~\cite{rafailov2023direct} are reported as
  small-scale single-objective comparators in
  Section~\ref{exp3-results}; together they bracket the preference-data
  spectrum, KTO needing only binary feedback and DPO needing pairwise.

  HITL-3 supplies two parallel feedback streams. \emph{Quality}
  feedback (helpfulness, specificity, empathy, accuracy, actionability)
  is the reward; \emph{compliance} feedback (promise safety, leakage,
  tone, legal) is the constraint. Feedback also propagates backward:
  classification corrections retrain Stage~1, rubric scores refine
  Stage~3 prompts, and response scores drive the RLHF updates, closing
  the Stage~5$\to$Stage~1 loop in Fig.~\ref{fig:architecture}.
%% =====================================================================
\section{Experimental Setup}\label{experiments}

\subsection{Dataset, Deduplication, and Annotation Provenance}
We use the RRGen corpus~\cite{gao2019automating}: 310{,}031 review-response pairs from 58 Android applications. After exact-text deduplication that removes $94{,}448$ duplicate reviews ($-30.5\%$ of the raw corpus), 215{,}583 unique reviews form the working corpus, labeled across the seven Maalej~\cite{maalej2015bug} classes. Evaluation additionally uses a 490-review expert gold standard (70 per class $\times$ 7 classes) and 100 stratified clusters for Stage~3 evaluation.

\textbf{Annotation provenance (full corpus, 215{,}583 reviews).} Of the working corpus, 79.49\% of labels come from the V2 LLM annotator with no override, 18.08\% are cleanlab-corrected against the verified anchor, and 2.43\% are directly human-verified by the lead author. The lead-author-versus-V2-LLM deviation is 25\%, replicated independently across two strata: the 5{,}041-review praise stratum and the 5{,}230-row human-verified anchor subset (Table~\ref{tab:annot-deviation}). This $\approx 25\%$ deviation rate, an order of magnitude higher than i.i.d.\ label-noise assumptions, motivates the confident-learning procedure of Stage~1.

\subsection{Annotation Methodology}
Five kinds of annotator contribute, following the multi-annotator convention of Gilardi et al.~\cite{gilardi2023chatgpt} and Pangakis et al.~\cite{pangakis2023automated} who treat LLM annotators as independent raters for inter-rater statistics.

\textbf{(i)~Lead author (domain expert).} Verified 5{,}230 reviews across minority strata (approximately thirty person-hours), wrote 20 reference IssueSpecs, labeled the 490-review classification gold blinded to V2 LLM, and scored the 400-row Stage~4 set on quality (1--5), specificity (1--5), and helpful (Y/N) with method IDs hidden until aggregation.

\textbf{Label provenance, stated plainly.} No LLM-generated label is used as ground truth
anywhere in this paper. LLM labels are the \emph{object under correction}, not the yardstick.
The V2 corpus labels are treated as noisy input, the confident-learning step proposes
corrections against a human-verified anchor, and every reported accuracy or $\kappa$ is measured
against a gold standard produced by human annotators. Where LLM raters do appear as raters, in
the 99-review cluster-purity check, they are labelled as such and no headline claim rests on
that panel.

\textbf{How rater conflicts are resolved.} The 490-review gold is the majority vote of three
independent human annotators over the seven-class label. Of the 489 usable items, 476 (97.3\%)
are unanimous and 13 (2.7\%) are two-to-one splits resolved by the majority; there are no
three-way ties, so no item required a tie-break rule or a further adjudication round. One item
was marked wrong by all three raters without a replacement label supplied and is excluded rather
than guessed. Disagreement is not discarded silently: the per-item votes are released with the
gold standard, so a reader can recompute any statistic on the unanimous subset alone. Before the
main task all three annotators completed a shared 20-review calibration round whose
disagreements were discussed and resolved against the written decision rules, which is where the
systematic conflicts (slow versus broken, device-specific versus general) were settled rather
than in the main pass.

\textbf{(ii)~Independent human rater panel (three raters).} The 490-review classification gold is annotated independently by the lead author and two further human raters on the identical item set. Each rater sees the classifier's predicted label and marks it correct (Y) or wrong (N), supplying the correct class when N. All three complete a shared 20-review calibration round with adjudicated disagreements before the main task opens, and work from the same written decision rules (Appendix~\ref{appendix:annotators}). The gold label for each review is the majority vote over the three raters. This is a verification task rather than open labeling, which is why agreement is high relative to free-labeling studies of app reviews.

\textbf{(iii)~Stage~4 rater panel (same three people).} Stage~4 was rated twice. The first round scored a 400-row set whose conditions were produced by template composers; we discarded it, for reasons given in Section~\ref{exp2-results}. The round we report is a 200-row set in which both arms come from the same model and differ only in whether the IssueSpec was supplied. All three raters scored all 200 rows independently, blind to the condition behind each reply, and we release both rounds together with the independence check that separates them. Three human annotators contribute in total, the lead author and two collaborators; categories (iv) and (v) are model raters, not people.

\textbf{(iv)~LLM-as-judge (Qwen2.5-3B-Instruct).} Scored 28 IssueSpecs on the five-dimension rubric and 50 sub-clusters on the Y/P/N purity rubric~\cite{zheng2023judging}, both calibrated against the lead author.

\textbf{(v)~Corpus-labeling LLMs.} The V2, cleanlab-V2, and V5 LLMs are treated as independent raters for inter-rater statistics.

Table~\ref{tab:annot-deviation} consolidates agreement and disagreement across all pairs.

\begin{table}[t]\centering\footnotesize
  \caption{Annotator-vs-LLM and inter-annotator deviation. The $\approx 25\%$ deviation is
  replicated on two independent strata; the $\kappa$-progression reaches moderate
  agreement~\cite{landis1977measurement}; the V5 endorsement row is a consistency check, not independent validation, since V5 trained on those corrections;
  LLM-judge calibration $\Delta$ bounds judge bias.}\label{tab:annot-deviation}
  \begin{tabular}{@{}p{4.0cm}p{1.5cm}p{2.5cm}@{}}\toprule
  Annotator pair & sample & agreement / deviation \\\midrule
  Lead-author vs V2 LLM & 10{,}271 & 75\% / \textbf{25\%} mislabeled \\
  Lead vs V2 / cleanlab / V5 & 490 gold & $\kappa$: $0.163 \to 0.333 \to \mathbf{0.592}$ \\
  3 human raters, Y/N verdict & 490 gold & Fleiss $\kappa = \mathbf{0.968}$; 479/490 unanimous \\
  3 human raters, 7-class label & 489 gold & Fleiss $\kappa = \mathbf{0.979}$; $\alpha = 0.979$ \\
  V2 LLM vs V5; cleanlab vs V5 & 215{,}583 & 71.96\% / 86.77\% agree \\
  cleanlab vs V5 (consistency, not independent) & 40{,}291 & 88.66\% support \\
  V5 vs gold, strictly held out & 307 & $\kappa = \mathbf{0.616}$ \\
  3 human raters (Stage~4 quality) & 400 rated rows & weighted $\kappa$ 0.970--0.986; within-1 100\% \\
  3 human raters (Stage~4 helpful) & 400 rated rows & $\kappa$ 0.678--0.836; $\alpha = 0.781$ \\
  3-rater \emph{LLM-assisted} purity check (lead + 2 Qwen prompts) & 99 reviews & Krippendorff $\alpha = 0.451$ \\
  Lead vs Qwen-judge (purity, rubric) & 50 / 28 & $\Delta = -0.035$ / $+0.23$ \\
  \bottomrule\end{tabular}\end{table}

Ratings are persisted as spreadsheets with JSON master keys, blinding shuffles conditions per pair, and the assignment unseals only at aggregation. Two distinct three-rater statistics appear in Table~\ref{tab:annot-deviation} and should not be conflated. The classification gold uses three independent \emph{human} raters on 490 reviews and reaches Fleiss $\kappa = 0.968$ on the Y/N verdict and $0.979$ on the effective seven-class label, with no three-way ties. The $\alpha = 0.451$ row is a separate cluster-purity check on 99 reviews that uses the lead author plus two Qwen2.5-3B prompts (concise and chain-of-thought) as raters.

\subsection{Multi-LLM Bias Control}\label{multi-llm-setup}
 A single-LLM headline raises a structural-bias concern: results from any one model may not
  transfer across families. We therefore replicate Stage~3 across four LLMs drawn from three
  distinct families, Anthropic Claude (Opus 4.7), Meta Llama-3.3-70B-Instruct (via Groq), and
  Alibaba Qwen2.5-\{3B, 1.5B\}-Instruct, with Stage~4 replicated on a Claude-vs-Qwen-3B subset,
  and check that cross-family rank order is preserved (Table~\ref{tab:multi-llm}).
  Frontier-family replication on GPT-4o and Gemini-1.5-Pro is queued.



\subsection{Experimental Design}
We pre-register three experiments, one per research question:

\textbf{Experiment~1 (RQ1, translation quality).} Five conditions are compared on the same 100 stratified clusters: (a)~LLM with taxonomy grounding (the full Stage~3 system), (b)~LLM without taxonomy grounding (free-form), (c)~raw review summary (a lower bound that concatenates the cluster's reviews), (d)~lead-author reference, and (e)~human-written GitHub issues mined from three open-source Android projects (the practitioner-quality upper bound). Each spec is scored on the five-dimension rubric. Significance is reported via paired Wilcoxon
signed-rank tests on rubric scores between conditions.

\textbf{Experiment~2 (RQ2, coupled vs.\ uncoupled).} The same 100 clusters drive four response-generation conditions: (a)~an rrgen-style baseline (review only), (b)~a prompt-baseline (CoRe-style with dev-rel prompt and lightweight entity extraction), (c)~RAG without IssueSpec, and (d)~the full system with IssueSpec. Evaluation uses automatic metrics (BLEU, ROUGE-L, BERTScore) and 400 blinded paired human ratings (quality 1--5, specificity 1--5, helpful Y/N), with the lead author as primary annotator and a second rater on a 30-pair subset. Statistical significance is checked simultaneously by paired Wilcoxon, Friedman with Nemenyi post-hoc~\cite{demvsar2006statistical}, Bradley-Terry, and McNemar at Bonferroni-corrected $\alpha = 0.05$. Effect sizes follow Romano et al.~\cite{romano2006appropriate} and Cliff~\cite{cliff1993dominance}.

\textbf{Why these four conditions.} The Experiment~2 ladder is designed so each rung adds
exactly one thing. The prompt baseline adds dev-rel system instructions on top of an rrgen-style
review-only baseline, \texttt{no\_spec} then adds retrieval, and \texttt{full} then adds the
validated IssueSpec. Reading the ladder from bottom to top isolates which single addition moves
human-rated quality. The headline comparison runs single-shot vanilla RAG~\cite{lewis2020retrieval}
rather than an agentic loop, so that the IssueSpec's contribution is not confounded with
orchestration; the agentic-RAG comparison against Self-RAG~\cite{asai2024self} and
DSP~\cite{khattab2022demonstrate} runs separately on a 10-review feasibility sample (stratified,
two per issue type, Qwen2.5-3B as critic and reviser, up to two iterations) and is reported as a
demonstration rather than a powered comparison.

\textbf{Experiment~3 (RQ3, dual- vs.\ single-objective RLHF).} Five policies run on a
  distilGPT2 proof-of-concept (82M parameters, 400 SFT samples, 296 KTO pairs, 100 DPO pairs):
  SFT-base, single-objective KTO~\cite{ethayarajh2024kto}, single-objective
  DPO~\cite{rafailov2023direct}, a constrained-proxy variant, and dual-objective Lagrangian
  Constrained PPO~\cite{schulman2017proximal, dai2024safe}. We measure head-to-head BLEU, ROUGE-L, and
  BERTScore on a 100-review test set, plus the rule-based safety-violation rate.
%% =====================================================================
\section{Results}\label{results}

\textbf{How to read this section.} Each research question owns one subsection, one headline
figure, and one headline table, and nothing else in the section is required to follow the
argument. Section~\ref{exp1-results} answers RQ1, is the typed IR producible at scale, with
Table~\ref{tab:stage3} for template-fill and Figure~\ref{fig:kappa} for the classifier recovery
that makes corpus-scale production possible. Section~\ref{exp2-results} answers RQ2, is the IR
load-bearing downstream, with Table~\ref{tab:stage4-rerun}; the answer is no, and the subsection explains what our earlier positive answer was measuring.
Section~\ref{exp3-results} answers RQ3, does constrained alignment demonstrate a
quality-versus-compliance trade-off, and reports a negative result.
Sections~\ref{cluster-validation-results} and~\ref{ablations} are supporting evidence: they test
whether individual components earn their place, and Section~\ref{load-bearing} states plainly
which ones do. A reader tracking only the contribution can read
Sections~\ref{exp1-results}, \ref{exp2-results}, and~\ref{load-bearing}.

Table~\ref{tab:findings} states what we take away from the section before the evidence for it,
so a reader can decide which subsections to read.

\begin{table}[t]\centering\footnotesize
\caption{What the experiments establish, in one line each. Negative findings are marked, and
they are findings rather than caveats.}\label{tab:findings}
\begin{tabular}{@{}p{0.6cm}p{7.6cm}p{2.4cm}@{}}\toprule
& finding & evidence \\\midrule
F1 & Type routing, not the language model, is what produces tracker-ready fields: removing it collapses strict fill from 0.96 to 0.69 and the per-type structural metrics to zero. & \S\ref{exp1-results}, A3 \\
F2 & LLM annotation noise on this corpus is structurally biased toward minority classes, not spread uniformly, which is why correcting against a small verified anchor works at all. & \S\ref{exp1-results} \\
F3 & \emph{Negative, and it overturns our own earlier claim.} With the generator held fixed, adding the typed IR changes response quality by $+0.03$ ($p=0.38$); the $+2.35$ we previously reported was an artifact of comparing two different text composers. & \S\ref{exp2-results} \\
F4 & Generator capability shows up mostly as whether fields get filled, not as the quality of the text once filled: template-fill spans 0.96--0.48 across four models while rubric quality spans 4.09--3.77. & \S\ref{exp1-results} \\
F5 & \emph{Negative.} The knowledge graph loses to count-matched flat clustering on all three intrinsic metrics; its defensible role is centrality-based prioritization. & \S\ref{cluster-validation-results}, A1b \\
F6 & \emph{Negative.} Extractive coverage does not measure faithfulness: \textsc{SpecCov} shows no rank correlation with human judgement and orders generators backwards. & \S\ref{speccov-validation} \\
F7 & \emph{Negative, then partial.} A CMDP constraint cannot be tested on a policy incapable of violating it; with one that can, the constraint binds and $\lambda$ grows, but no trade-off is demonstrated. & \S\ref{exp3-results} \\
F8 & Run-to-run generation noise is about 0.03 on template-fill, which is the yardstick separating the differences we claim from the ones we do not. & \S\ref{reproducibility} \\
\bottomrule\end{tabular}\end{table}

\subsection{Is the IR Producible at Scale?}\label{exp1-results}

Under strict content-validity criteria (at least three reproduction steps with action verbs, descriptions $\geq 30$ words, and a full As-a/I-want/So-that triple for user stories), the LLM-with-taxonomy condition achieves a substantive template-fill rate of 0.96 (Table~\ref{tab:stage3}). Free-form LLM reaches 0.69, the same as the lead-author reference set ($n=20$). Raw concatenation lands at 0.34, and real GitHub issues mined from three open-source Android projects reach 0.53. The LLM-vs-GitHub gap reflects coverage rather than content quality: the LLM is prompted to fill every field, whereas humans leave optional fields blank. A very-strict sensitivity sweep preserves the ordering, Claude drops from 0.97 to 0.71, with the Claude-vs-GitHub gap shifting only $-0.085$.

\textbf{\textsc{SpecCov} scores, and a correction to how they were computed.} An earlier
version of our scorer applied per-condition floors, raising \texttt{raw\_summary} to 5 and the
human reference to 4 regardless of content, on the reasoning that verbatim copying is faithful by
construction and hand-written references deserve a floor. That is not a measurement, because a
score whose value depends on the label of the thing being scored cannot compare labels, so we
removed the floors and report what the scorer produces unaided. Across 320 specs the
LLM-with-taxonomy condition averages 4.19, free-form LLM 3.38, raw concatenation 4.47, and the
human reference 3.45. The two LLM conditions are unchanged, because no floor was ever applied to
them; the floors had inflated the two comparison conditions, so removing them makes our own
system look better rather than worse, and we note that explicitly so the direction of the
correction is on the record. Read on their own these numbers would say the taxonomy condition
scores highest. Section~\ref{speccov-validation} shows they should not be read on their own. Top-ranked aspects by PageRank centrality are
\texttt{ad} (0.040), \texttt{phone} (0.011), and \texttt{battery} (0.007).

\subsubsection*{SpecCov does not measure faithfulness}\label{speccov-validation}
An extractive-coverage score is only useful if it agrees with what a person calls faithful. We
tested that directly. Three human raters independently scored 60 specs, 20 from each of the
three generators, shuffled and unlabelled, on a 1--5 faithfulness scale, and separately listed
every concrete claim the source reviews do not support. One qualification is needed before the
numbers. One rater completed the corrected sheet, which shows exactly the evidence
\textsc{SpecCov} itself scores against; the other two completed an earlier draft that showed a
subset of it, so they judged against less evidence than the metric uses and their flags are
biased toward over-flagging. We therefore lead with the rater who saw the full evidence and
report the other two beside her. The metric fails all three checks on either reading.

\emph{It does not rank specs the way humans do.} Spearman $\rho$ between \textsc{SpecCov} and
human faithfulness is $+0.045$ ($p = 0.73$), $+0.015$ ($p = 0.91$), and $+0.226$ ($p = 0.08$) for
the three raters, and $+0.147$ ($p = 0.26$) against their mean. No correlation reaches
significance.

\emph{It does not detect the claims humans call unsupported.} On the 60 specs, one rater marked
45 as containing at least one unsupported detail and 15 as clean. \textsc{SpecCov} scores those
two groups at 3.64 and 3.93, a difference in the right direction but far from significant
(Mann-Whitney $p = 0.34$). A metric that detected hallucination would separate these groups.

\emph{It orders the generators backwards.} \textsc{SpecCov} ranks the taxonomy condition highest
(4.20) and the human reference lowest (3.45); the humans rank them the other way, human
reference 4.10 and taxonomy 3.50. The mechanism is visible in the definition: templated specs
reuse source vocabulary densely, which is what the score rewards, while a human writer condenses
and paraphrases, which the score penalizes even when nothing is invented.

We therefore withdraw \textsc{SpecCov} as a faithfulness measure. Lexical grounding is not a
usable proxy for faithfulness on this task, and the automatic 320-spec rubric that appeared to
support it is not independent evidence, since its faithfulness dimension is computed by the same
overlap procedure. Entailment-based measurement~\cite{laban2022summac, honovich2022true,
durmus2020feqa} is the direction we would now take. We release the scorer, the 60-spec human
ratings, and this analysis, because a documented negative result on a plausible-looking metric
is more useful to the next person than a metric that was never tested.

\textbf{Cross-LLM replication (bias control).} Table~\ref{tab:multi-llm} reports four uses across three model families. Llama-3.3-70B reaches 0.80 template-fill, Qwen2.5-3B 0.74, and Qwen2.5-1.5B 0.48 under the same strict criteria; the templated~$>$~free-form~$>$~raw ordering holds across all four models with capability-scaled magnitudes. On two per-class structural-compliance checks, Llama-3.3-70B attains complete template compliance on the stratified sample: all 29 bug specs include three or more reproduction steps with action verbs (vs.\ Claude's 73\%), and all 30 feature specs include the As-a/I-want/So-that triple (vs.\ Claude's 97\%). These measure structural template compliance (does the field exist with the required form?), not content correctness, content quality is reported separately under the five-dimension rubric. The structural-bias control that matters here is cross-family rank preservation: the IssueSpec contract transports across families rather than depending on a single proprietary generator.

\textbf{Cross-family rubric quality.} Template-fill measures whether a generator emits the
required fields; it says nothing about whether the content is any good. We therefore ran the same
Qwen2.5-3B judge and the same five-dimension rubric over the other generators' specs, on matched
clusters (Table~\ref{tab:xfam-rubric}). On the 14 clusters all four generators cover, mean rubric
score follows the same order as template-fill: Claude 4.09, Llama-3.3-70B 3.96, Qwen2.5-3B 3.83,
Qwen2.5-1.5B 3.77. We state plainly what that ordering is worth. At $n=14$ none of the six
pairwise differences reaches significance (Wilcoxon $p$ between 0.08 and 0.50), and the adjacent
gaps are the same size as the run-to-run generation noise measured in
Section~\ref{reproducibility}. The four-way ordering is therefore consistent with the
template-fill result rather than independent confirmation of it. The Claude-versus-Llama
comparison in Panel B, which has twice the sample and spans all five issue types, does reach
significance at $+0.38$ ($p = 0.015$), and that is the one cross-family quality difference we
assert. The spread is much narrower than the template-fill spread (0.96 to 0.48),
which is informative in itself: the capability gap between these models shows up mostly as
whether fields get filled at all, not as the quality of the text once they are. Completeness is
the dimension that separates them (4.00 down to 3.00), while accuracy is nearly flat (4.50 for
three of the four). On the wider Claude-versus-Llama panel, which spans all five issue types,
Claude leads 4.01 to 3.64, with the gap concentrated in actionability and completeness. Two
caveats belong with these numbers. The four-way panel is bug-report clusters only, because that
is where the two Qwen runs overlap, and the judge is itself Qwen2.5-3B-Instruct, so the
Qwen2.5-3B row carries a self-preference risk that the others do not.

\begin{table}[t]\centering\footnotesize
\caption{Cross-family five-dimension rubric, same Qwen2.5-3B judge as
Table~\ref{tab:rubric}, on matched clusters. Panel A is the four-way comparison on clusters
covered by every generator (bug reports only). Panel B widens the Claude-versus-Llama comparison
to all five issue types. The Panel-A ordering matches the template-fill ordering in
Table~\ref{tab:multi-llm}, but none of its six pairwise differences is significant at $n=14$
(Wilcoxon $p$ from 0.08 to 0.50), so read it as descriptive. Panel B, with twice the sample and
all five issue types, does separate: $+0.38$ for Claude over Llama, Wilcoxon
$p = 0.015$.}\label{tab:xfam-rubric}
\begin{tabular}{@{}lrrrrrrr@{}}\toprule
generator & $n$ & mean & compl. & accur. & action. & spec. & clarity \\\midrule
\multicolumn{8}{@{}l}{\emph{Panel A, four-way, 14 matched bug-report clusters}} \\
Claude (Anthropic) & 14 & \textbf{4.09} & 4.00 & 4.50 & 3.71 & 4.36 & 3.86 \\
Llama-3.3-70B (Meta) & 14 & 3.96 & 3.86 & 4.50 & 4.07 & 3.86 & 3.50 \\
Qwen2.5-3B (Alibaba) & 14 & 3.83 & 3.07 & 4.50 & 4.00 & 3.86 & 3.71 \\
Qwen2.5-1.5B (Alibaba) & 13 & 3.77 & 3.00 & 4.15 & 4.46 & 3.62 & 3.62 \\
\midrule
\multicolumn{8}{@{}l}{\emph{Panel B, Claude vs Llama, 28 clusters, all five issue types}} \\
Claude (Anthropic) & 28 & \textbf{4.01} & 3.86 & 4.54 & 3.75 & 4.29 & 3.64 \\
Llama-3.3-70B (Meta) & 28 & 3.64 & 3.39 & 4.46 & 3.14 & 3.82 & 3.36 \\
\bottomrule\end{tabular}\end{table}

\begin{table}[t]\centering\footnotesize
\caption{Multi-LLM evidence summary across three model families (Anthropic, Meta, Alibaba). Cross-family rank preservation across Stage~3 and Stage~4 is the structural-bias control; LLM-as-judge calibration $\Delta$ vs the lead author bounds judge bias.}\label{tab:multi-llm}
\begin{tabular}{@{}p{2.7cm}p{2.9cm}p{2.4cm}@{}}\toprule
Use & Model (family) & Result \\\midrule
Stage~3 generate & Claude (Anthropic) / Llama-3.3-70B (Meta) / Qwen-3B (Alibaba) / Qwen-1.5B (Alibaba) & 0.96 / 0.80 / 0.74 / 0.48 \\
Stage~4 generate & Claude / Qwen-3B & 0.61 / 0.54 \\
5-dim rubric judge & Qwen-3B, all 100 specs & 3.94/5 (28-spec subsample: 3.89) \\
Cross-family rubric & Claude / Llama / Qwen-3B / Qwen-1.5B & 4.09 / 3.96 / 3.83 / 3.77 (Table~\ref{tab:xfam-rubric}) \\
Y/P/N purity judge & Qwen-3B vs lead-author & 0.625 vs 0.66; $\Delta=-0.035$ \\
\bottomrule\end{tabular}\end{table}

\textbf{Five-Dimension Rubric.} The proposal specifies a 5-dimension Likert rubric
(completeness, accuracy, actionability, specificity, clarity, each 1--5)~\cite{gilardi2023chatgpt,
pangakis2023automated, zheng2023judging}, operationalized with Qwen2.5-3B-Instruct as the judge.
We score the complete set rather than a subsample: all 100 taxonomy specs and all 20 lead-author
reference specs, in a single run (Table~\ref{tab:rubric}). The taxonomy condition reaches
$3.94/5$ (sd 0.59 across specs), above the hypothesized lower bound of 3.5. An earlier stratified
28-spec subsample under the same judge and rubric gave $3.89$, so subsampling cost $0.05$ of a
Likert point here, which bounds how much the small-sample version of this number could have been
misleading. Per issue type the taxonomy condition ranges from 4.23 on performance and 4.02 on
bug reports down to 3.45 on usability, the type whose reviews are least specific about what
went wrong.

The reference set is a judge calibration, not a competitor. It scores $3.45$ overall, below the
generated specs, and the gap is concentrated in actionability (2.85 versus 3.86). The reason is
visible in the sheets: the lead author filled title, description, severity, and component but
left the template-specific fields such as reproduction steps and acceptance criteria empty,
because writing a reference spec by hand is a different task from filling a routed template. The
comparison therefore says that a human writing freely produces less tracker-ready output than a
type-routed generator, which is the same effect the 0.96-versus-0.53 template-fill gap measures
against real GitHub issues. It is not evidence that the specs are better than expert judgement
on content, and we do not read it that way.

\begin{table}[t]\centering\footnotesize
\caption{Experiment~1 five-dimension rubric (Qwen-judge, 1--5 scale) on the \emph{complete}
spec set: all 100 taxonomy specs and all 20 lead-author reference specs, scored in one run so
the two columns are directly comparable. The middle column repeats the earlier stratified
28-spec subsample under the same judge and rubric; the gap between it and the full set is
$+0.05$, which is the subsample error. The LLM-as-judge protocol follows the
calibration-against-reference convention of~\cite{gilardi2023chatgpt, pangakis2023automated,
zheng2023judging}.}\label{tab:rubric}
\begin{tabular}{lrrr}\toprule
dimension & taxonomy, $n=100$ & taxonomy, $n=28$ & lead-author ref, $n=20$ \\\midrule
completeness & \textbf{3.72} & 3.79 & 3.20 \\
accuracy & \textbf{4.43} & 4.39 & 4.10 \\
actionability & \textbf{3.86} & 3.71 & 2.85 \\
specificity & \textbf{4.15} & 4.07 & 4.05 \\
clarity & \textbf{3.52} & 3.50 & 3.05 \\\midrule
overall mean & \textbf{3.94} & 3.89 & 3.45 \\
sd across specs & 0.59 & -- & 0.41 \\
\bottomrule\end{tabular}\end{table}

\begin{table}[t]\centering\footnotesize
  \caption{Experiment~1 substantive template-fill (strict criteria). The 0.96 vs 0.53
  LLM-vs-GitHub gap is coverage, not content quality; the strict-vs-very-strict sweep preserves
  ordering. The exact-zero cells in \texttt{bugs:repro} and \texttt{features:story} columns for
  (b), (c), (d), and (e) are structural: free-form LLM, raw concatenation, lead-author, and
  human-GitHub specs do not emit the template-specific reproduction-step list or
  As-a/I-want/So-that triple at all when not prompted for them, so the strict-fill rate is zero
  by construction.}\label{tab:stage3}
  \begin{tabular}{lrrr}
  \toprule
  condition & strict & bugs:repro & features:story \\
  \midrule
  (a) LLM + taxonomy & \textbf{0.96} & \textbf{0.73} & \textbf{0.97} \\
  (b) LLM free-form & 0.69 & 0.00 & 0.00 \\
  (c) raw concat & 0.34 & 0.00 & 0.00 \\
  (d) lead-author ref ($n=20$) & 0.69 & 0.00 & 0.00 \\
  (e) human GitHub (3 repos) & 0.53 & 0.13 & 0.00 \\
  \bottomrule
  \end{tabular}
  \end{table}

\textbf{Stage~1 Classification Recovery.} The Cohen's $\kappa$ trajectory against the 490-review expert gold crosses three Landis-Koch~\cite{landis1977measurement} agreement bands (Table~\ref{tab:kappa}): the V2 LLM achieves $\kappa = 0.163$ (slight); after cleanlab correction~\cite{northcutt2021confident} $\kappa = 0.333$ (fair); the V5 classifier retrained on the corrected labels reaches $\kappa = 0.592$ (moderate, just below the substantial threshold). The trajectory does not depend on any single annotator. Scored against the majority-vote gold of the three-rater human panel ($n = 489$ after excluding one review left without a correction label), the same three classifiers reach $\kappa = 0.165$, $0.334$, and $0.590$, within $0.002$ of the lead-author values at every step (Table~\ref{tab:kappa}). On the strictly held-out 307-review subset (183 of the 490 entered training via the verified anchor), V5 reaches $\kappa = 0.616$, slightly above the 0.592 on the full 490, ruling out a contamination artifact.

The per-class recovery story (Table~\ref{tab:perclass-f1}) makes the structural collapse visible. On the two minority classes the V2 LLM cannot see at all, F1 jumps to 0.83 for \texttt{compatibility} and 0.77 for \texttt{performance}. Applied across the full 215{,}583-review corpus, V5 agrees with the cleanlab correction on
88.66\% of the 40{,}291 corrected cases. We want to be exact about what that number is and is
not, because it is easy to over-read. V5 was trained on a corpus in which 38{,}985 of those
corrected rows carry the corrected label, so its agreement is substantially a measure of fit to
its own training signal, not an independent second opinion. We report it as a consistency check,
which is what it is, and we do not claim the two signals are independent.

The stronger evidence is the held-out gold, and we describe the split precisely because the
mechanism matters. V5's training corpus is built with a stratified per-class cap of 15{,}000
rows at seed 42. Replaying that cap identifies which corpus rows survived into training: 183 of
the 490 gold reviews did, and 307 were capped out and never seen. The split is therefore an
artifact of class-balancing, not of anchor membership, and it is unbiased with respect to
difficulty for exactly that reason. On the 307 unseen reviews V5 reaches $\kappa = 0.616$,
slightly \emph{above} the 0.592 it scores on the full 490. A model that had merely memorised its
training labels would fall on unseen items rather than hold. One caveat belongs with this. The
split is reconstructed by replaying the cap rather than read from a saved training manifest, so
it is a faithful reconstruction rather than a recorded fact; we release the replay script so the
reconstruction can be checked. The choice of confident-learning thresholds is not load-bearing: a full $3\times 3$ grid over $(\text{anchor\_conf}, \text{max\_llm\_prob}) \in \{0.65, 0.70, 0.75\}\times\{0.15, 0.20, 0.25\}$ keeps V5 endorsement in an 82.7--91.0\% band, with the headline $(0.70, 0.20)$ at 88.66\% sitting near the grid center.

\begin{figure}[t]\centering
  \includegraphics[width=\linewidth]{fig_kappa_progression.png}
  \caption{Stage-1 classifier recovery and the gold standard it is measured against. \emph{Left:}
  Cohen's $\kappa$ for the three classifier versions, plotted twice, once against the
  lead-author gold (solid) and once against the three-rater majority-vote gold (hatched). The
  two are within 0.002 at every step, so the 0.163 to 0.592 trajectory does not depend on any
  single annotator. Background bands are the Landis-Koch~\cite{landis1977measurement} agreement
  ranges. \emph{Right:} agreement among the three human raters on the same 490 reviews, pairwise
  and as a three-rater Fleiss $\kappa$, all above the almost-perfect threshold. Read the left
  panel for the result and the right panel for the reliability of the yardstick it uses.}
  \label{fig:kappa}
\end{figure}

\begin{table}[t]\centering\footnotesize
  \caption{Cohen's $\kappa$ progression vs the 490-review expert gold. Three
  Landis-Koch~\cite{landis1977measurement} bands crossed; cleanlab correction follows the
  confident-learning procedure of~\cite{northcutt2021confident}. The final column repeats the
  measurement against the majority-vote gold of the three-rater human panel, showing the
  trajectory is not an artifact of a single annotator.}\label{tab:kappa}
  \begin{tabular}{lrrlr}\toprule
  classifier & $n$ & Cohen's $\kappa$ & Landis-Koch band & $\kappa$ vs 3-rater gold \\\midrule
  V2 LLM original & 490 & 0.163 & slight & 0.165 \\
  cleanlab-corrected V2 & 490 & 0.333 & fair & 0.334 \\
  \textbf{V5 classifier} & 490 & \textbf{0.592} & \textbf{moderate} & \textbf{0.590} \\
  \bottomrule\end{tabular}\end{table}

\begin{table}[t]\centering\footnotesize
  \caption{Per-class F1 against the 490-review expert gold. V2 LLM is blind to the two minority
  classes (compatibility, performance). The verified-anchor procedure with targeted augmentation
  recovers them.}\label{tab:perclass-f1}
  \begin{tabular}{lrrr}\toprule
  class & V2 LLM & cleanlab-V2 & \textbf{V5} \\\midrule
  \texttt{bug\_report} & 0.43 & 0.50 & \textbf{0.71} \\
  \texttt{feature\_request} & 0.59 & 0.61 & \textbf{0.69} \\
  \texttt{performance} & 0.00 & 0.47 & \textbf{0.77} \\
  \texttt{usability} & 0.11 & 0.21 & \textbf{0.55} \\
  \texttt{compatibility} & 0.00 & 0.00 & \textbf{0.83} \\
  \texttt{praise} & 0.41 & 0.55 & \textbf{0.79} \\
  \texttt{other} & 0.27 & 0.36 & \textbf{0.41} \\\midrule
  macro F1 & 0.26 & 0.39 & \textbf{0.68} \\
  \bottomrule\end{tabular}\end{table}

To rule out memorization of the 200 synthetic compatibility templates used during V5 training, we evaluate V5 against Maalej's 5{,}008 labels under a label-space cross-protocol test. Review texts appear in the anchor under the RRGen seven-class taxonomy, but predictions are scored against Maalej's separate taxonomy, which contains no \texttt{compatibility} class, making any compatibility predictions out-of-distribution with respect to the evaluation labels. On the five overlapping classes ($n=4{,}560$), V5 achieves macro F1~=~0.676, weighted F1~=~0.730, and accuracy~=~72.0\%. V5 also flags 432 reviews (8.6\%) as \texttt{compatibility}; the top-twenty-confidence flags read as canonical compatibility complaints (``crashes on iphone 5 and ipad 3'', ``works on HTC Desire X but not Samsung Galaxy Tab 3''), all labeled \texttt{bug\_report} in Maalej because Maalej offers no alternative class. The detector is performing concept-recognition, not template-matching.

\subsection{Is the IR Load-Bearing? A Result We Had to Withdraw}\label{exp2-results}

An earlier version of this experiment reported that adding the IssueSpec to the response
generator raised human-rated quality by $2.35$ Likert points. That number was wrong, not in its
arithmetic but in what it measured, and the correction is the most important thing this section
has to report.

\textbf{What was wrong.} The four conditions in that experiment were not produced by a language
model. Each was produced by a deterministic template composer, and the composers differed in
authorial register as much as in what information they were given. The \texttt{no\_spec} arm
emitted lemmatised, ungrammatical text averaging 78 words (``we be sorry that the login loop you
describe have got in the way''); the \texttt{full} arm emitted polished prose averaging 123 words,
with cleanup passes and an empathy phrasebook that no other arm received. A rater comparing them
was separating two writing styles, and the presence or absence of the IssueSpec was confounded
with that separation. The gap measured the composer at least as much as the specification.

\textbf{The clean comparison.} We re-ran the experiment with the confound removed. Both arms are
now generated by the same model, given the same retrieved context, on the same 100 reviews, and
differ in exactly one respect: whether the validated IssueSpec appears in the prompt. The two
arms average 105 and 97 words with distinct-1 of 0.736 and 0.760, so they are not separable on
surface features. Three raters scored all 200 blinded rows.

\textbf{The IssueSpec does not measurably improve response quality.} Pooled over the three
raters the gain is $+0.03$ Likert points (position-controlled; $+0.04$ uncontrolled), 95\%
bootstrap CI $[-0.10, +0.18]$, paired Wilcoxon $p = 0.38$ on 100 pairs, Cliff's $\delta = 0.04$.
Specificity behaves the same way, pooled $-0.00$ ($p = 0.87$), and the helpful rate is flat at
94\% in both arms. On the evidence of this design we cannot claim that supplying the typed
specification improves the response a user receives.

\textbf{The raters disagree with each other, and that is the interesting part.}
Table~\ref{tab:stage4-rerun} breaks the result out. The lead author, who designed the system,
prefers the IssueSpec arm decisively: $+0.53$ quality ($p = 2 \times 10^{-11}$, 59 wins, 5
losses, 36 ties) and $+0.39$ specificity ($p < 10^{-7}$). The two independent raters lean the
other way, $-0.26$ and $-0.18$, neither significant at $\alpha = 0.05$ ($p = 0.06$ and $0.09$).
Inter-rater reliability is correspondingly low, quadratic-weighted $\kappa$ of 0.03, 0.04 and
0.50, against 0.97 in the discarded template-based round. High agreement there was a symptom of
the confound: when one arm is visibly broken, raters agree easily. Here they do not.

We looked for a content mechanism behind the lead author's preference and did not find one.
Correlating each rater's per-review preference against the spec-derived features that the
IssueSpec should introduce, named devices, reproduction steps, and component mentions, produces
no correlation that survives inspection for any rater; the largest is $\rho = 0.19$ ($p = 0.06$)
between the lead author's preference and reply length. So we can describe the disagreement but
not explain it. Two readings remain open. The lead author may be perceiving triage value that a
general reader has no reason to weigh, in which case the right evaluation asks developers rather
than readers. Or the preference may be expectancy on the part of the person who built the system.
This design cannot separate those, and we are not going to assert the flattering one.

\begin{table}[t]\centering\footnotesize
\caption{Stage-4 re-run, same model with and without the IssueSpec, 100 paired reviews, three
raters, all 200 rows scored by each. The pooled effect is null. The disagreement between the
lead author and the two independent raters is larger than the effect being
measured.}\label{tab:stage4-rerun}
\begin{tabular}{@{}lrrrrr@{}}\toprule
rater & no\_spec & full & gain & 95\% CI & Wilcoxon $p$ \\\midrule
lead author & 4.04 & 4.60 & $+0.53$ & $[+0.44, +0.68]$ & $2 \times 10^{-11}$ \\
independent rater 1 & 3.87 & 3.58 & $-0.26$ & $[-0.56, -0.01]$ & 0.06 \\
independent rater 2 & 3.57 & 3.42 & $-0.18$ & $[-0.33, +0.03]$ & 0.09 \\\midrule
\textbf{pooled} & 3.83 & 3.87 & $\mathbf{+0.03}$ & $[-0.10, +0.18]$ & \textbf{0.38} \\
\bottomrule\end{tabular}\end{table}

\textbf{Why we believe the raters this time.} The discarded round showed an arithmetic signature
that independent rating cannot produce: the two additional raters deviated from the lead author
on 16 and 19 of 400 rows, those deviation sets did not overlap on a single row, and every
deviation was exactly one point. In this round the same check passes. The raters differ from the
lead author on 133 and 145 of 200 rows, those sets overlap on 101 where chance predicts about 96,
and the deviations range from one to four points. We release the check with the data because a
reader should be able to apply it to our sheets rather than take our word for it.

\textbf{What survives.} The negative finding about retrieval survives and is unaffected by this
correction, since it compares retrieval against a no-context baseline rather than against the
IssueSpec. What does not survive is the claim that the typed IR carries downstream response
quality. Stage~3's contribution, producing tracker-ready structured records at corpus scale,
stands on its own evidence in Section~\ref{exp1-results}; it no longer has a demonstrated
downstream payoff in response generation.

\subsection{Aligning the IR-Aware Generator at Scale}\label{exp3-results}

  Five policies were trained on distilGPT2 (82M parameters, 400 SFT samples, 296 KTO pairs, 100
  DPO pairs): SFT-base, KTO~\cite{ethayarajh2024kto}, DPO~\cite{rafailov2023direct}, a constrained-proxy
  variant, and Lagrangian Constrained PPO~\cite{schulman2017proximal, dai2024safe}. Head-to-head on a
  100-review test set, per-policy BLEU-1~/~ROUGE-L~/~BERTScore~F1 are: SFT-base
  0.090~/~0.097~/~0.802; KTO 0.068~/~0.074~/~0.792; DPO 0.084~/~0.088~/~0.800;
  \textbf{constrained-proxy 0.137~/~0.131~/~0.806}; Lagrangian PPO 0.087~/~0.091~/~0.798. The
  constrained proxy leads three of four metrics (BLEU-1, ROUGE-L, BERTScore F1; the rule-based
  safety-violation rate is uninformative here because compliance is satisfied at
  initialization), with $+52\%$ BLEU-1 over SFT-base. KTO underperforms here, consistent with
  its small-data assumption breaking on a 296-pair training set.

  \textbf{The dual-objective claim is not demonstrated, and we report that as the result.} At
  distilGPT2 scale the compliance scorer is satisfied at initialization ($C = 0.94 \geq \tau$),
  so the Lagrange multiplier $\lambda$ collapses to zero and the update reduces to
  single-objective optimization. Two further experiments confirm this is a property of the
  regime rather than of one threshold choice. A re-run that replaces the permissive scorer with
  the four operational compliance checks and tightens $\tau$ from 0.50 to 0.90 records 0 binding
  steps out of 30 and a single violation across the whole run. A sweep over
  $\tau \in \{0.0, 0.5, 0.7, 0.9, 1.0\}$ leaves the pass rate at 100\% for four of the five
  policies at every threshold, with the first binding $\tau$ undefined for those four; only
  SFT-base ever dips, to 99\% at $\tau \geq 0.9$. A model of this size cannot plausibly violate
  the rubric, so the constraint has no work to do and the CMDP machinery is untested under
  binding conditions. The $+52\%$ BLEU-1 figure for the constrained proxy is therefore a
  single-objective result, not evidence of a quality-versus-compliance trade-off, and we do not
  claim it as one.

  \textbf{The constraint binds once the policy can violate it.} The obstacle is not parameter
  count, it is capability: distilGPT2 never produces the fluent developer-relations phrasing the
  rubric penalises, so it cannot violate and the constraint has nothing to hold back. We
  measured this directly. Under the same operational scorer, 5.2\% of the paper's own Stage-4
  generations trip a check, every one of them by over-promising, and a zero-shot probe of
  Qwen2.5-1.5B-Instruct in the same reply setting gives an indistinguishable 5.0\%. distilGPT2
  gives roughly zero. We therefore re-ran the Lagrangian loop with Qwen2.5-1.5B-Instruct as the
  policy, LoRA adapters over the attention projections, the frozen base as the KL reference, and
  a zero-tolerance threshold $\tau = 1.0$; the scorers, the dual-ascent rule, and the batch size
  are unchanged from the distilGPT2 run. The behaviour inverts. Over 90 steps the constraint binds
  on 13 of them (14.4\%), and $\lambda$ leaves zero, rises monotonically to 0.65, and never
  collapses back. A 30-step run under the same setup reaches $\lambda = 0.125$, so the multiplier
  scales with the horizon rather than saturating. Split by thirds, the trajectory shows the dual
  ascent responding to violations as they arrive: 1 binding step and $\lambda = 0.05$ in the
  first thirty, 7 and $\lambda = 0.40$ in the second, 5 and $\lambda = 0.65$ in the third. The
  qualitative claim that the CMDP degenerates to single-objective optimization is therefore a
  property of the proof-of-concept policy, not of the formulation.

  \textbf{What is and is not shown.} The two objectives visibly pull against each other. Across
  the run quality rises from 0.565 to 0.590 while compliance slips from 0.997 to 0.989, and
  $\lambda$ climbs in response, which is the dual mechanism doing exactly what it is meant to do
  when the policy drifts toward the reward at the constraint's expense. What we cannot yet show
  is the other half of the story, the multiplier growing large enough to pull the policy back
  into compliance. At a 5\% base violation rate and a batch of four, roughly one violating sample
  arrives every seven steps, and 90 steps of LoRA updates at this learning rate do not accumulate
  enough constraint pressure to reverse the drift. So the correct statement is narrower than a
  confirmed RQ3 and wider than the earlier null: the dual objective engages, the multiplier
  responds monotonically to violations, and the tension between the objectives is measurable,
  while closing the loop back to compliance needs a longer horizon or a higher base violation
  rate than a 1.5B instruct model supplies.

  What this layer contributes is the CMDP formulation instantiated over four
  domain-specific dev-rel constraints, an end-to-end trainer that runs, the threshold sweep that
  establishes the non-binding regime, the capability probe that explains it, the binding-constraint
  re-run, and the released testbed.

\subsection{Cluster Quality and Purity}\label{cluster-validation-results}

\begin{figure}[t]\centering
  \includegraphics[width=\linewidth]{fig_cluster_purity.png}
  \caption{Extrinsic cluster purity for the flat baseline and the KG hierarchy, from the saved
  judge audit. \emph{Left:} weighted purity per issue type, flat-194 against KG-605. The KG loses
  on four of five types, which is the finer-granularity cost we report rather than a result in
  its favour. \emph{Right:} the underlying Y/P/N verdict counts behind the two overall numbers,
  0.625 for flat and 0.390 for the KG. This figure is the evidence for the A2 ablation and for
  the claim in Section~\ref{load-bearing} that the KG's contribution is prioritization, not
  cluster quality.}
  \label{fig:purity}
\end{figure}

  Stage~2 was redesigned from a flat UMAP+HDBSCAN baseline (194 clusters, mean size 51.5) to the
   three-layer KG-hierarchical design (605 sub-clusters, mean size 3.0). We compute two
  complementary validity axes. \emph{Intrinsic geometry},
  Davies-Bouldin~\cite{davies1979cluster}, Calinski-Harabasz~\cite{calinski1974dendrite},
  silhouette~\cite{rousseeuw1987silhouettes}, follows standard IR/clustering
  practice~\cite{schutze2008introduction, steinbach2000comparison, strehl2002cluster}.
  \emph{Extrinsic cluster purity}~\cite{amigo2009comparison, schutze2008introduction} uses a
  Y/P/N weighted-purity rubric, $\mathrm{purity}_w = (|Y| + 0.5|P|)/(|Y|+|P|+|N|)$, with Y~=~all
   5 share theme, P~=~3--4 of 5 share, N~=~fewer than 3 share, following inter-rater conventions
   compatible with Landis-Koch~\cite{landis1977measurement}.

  KG-hierarchical achieves $5.4\times$ lower Davies-Bouldin and $1.9\times$ higher
  Calinski-Harabasz than the flat baseline (Table~\ref{tab:cluster-quality}), at the cost of
  finer per-cluster Y/P/N purity (0.39 vs.\ 0.625 Qwen-judge). By Landis-Koch~\cite{landis1977measurement} convention adapted to weighted purity, the flat baseline's
  0.66 lead-author purity sits in the substantial-agreement band ($\geq 0.61$) and KG-605's 0.39
   in the fair band ($\geq 0.21$), reflecting the granularity-vs-coherence trade-off the A1b
  ablation makes explicit.

\begin{table}[t]\centering\footnotesize
  \caption{Stage~2 designs head-to-head, comparing Pure HDBSCAN (without UMAP), original flat
  UMAP+HDBSCAN, and KG-hierarchical. Intrinsic metrics (Davies-Bouldin~\cite{davies1979cluster},
   Calinski-Harabasz~\cite{calinski1974dendrite}, silhouette~\cite{rousseeuw1987silhouettes})
  follow standard IR/clustering convention~\cite{schutze2008introduction,
  steinbach2000comparison, strehl2002cluster}. The Y/P/N weighted-purity rubric is a
  discrete-coding instantiation of the extrinsic cluster-purity
  metric~\cite{amigo2009comparison, schutze2008introduction}. The count-controlled flat baseline
   (agglomerative-605) is reported separately in the A1b ablation            
  (\S\ref{ablations}).}\label{tab:cluster-quality}
  \begin{tabular}{@{}p{3.4cm}rrr@{}}\toprule
  metric & Pure HDB & Flat U+H & \textbf{KG hier.} \\\midrule
  n\_clusters & 3 (deg.) & 194 & \textbf{605} \\
  mean cluster size & $\sim$16{,}933 & 51.5 & \textbf{3.0} \\
  Davies-Bouldin ($\downarrow$) & n/a & 12.15 & \textbf{2.24} \\
  Calinski-Harabasz ($\uparrow$) & n/a & 0.98 & \textbf{1.85} \\
  Y/P/N purity (lead-author) & n/a & 0.66 & queued \\
  Y/P/N purity (Qwen-judge) & n/a & 0.625 & 0.390 \\
  \bottomrule\end{tabular}\end{table}
  
\textbf{Count-controlled diff from original.} A count-controlled ablation (A1b) disambiguates
  two evaluation axes that prior cluster-quality work tends to conflate: intrinsic geometry and
  practitioner-side prioritization (centrality-weighted coverage). On 1{,}815 representative
  reviews (3 per sub-cluster across 605 sub-clusters), an agglomerative clustering forced to 605
   clusters reaches DB = 1.12, CH = 4.21, silhouette = $+0.148$, against KG-605's 2.24 / 1.85 /
  $-0.234$. The $5.4\times$ DB advantage in the headline against flat-194 is therefore primarily
   a cluster-count effect, not a KG-structure effect: on geometry, count-matched agglomerative
  clustering ties or beats the KG.

  \textbf{Where the KG actually wins.} On the prioritization axis the story reverses. Of the
  10{,}534 aspect nodes in the KG, the top five by PageRank account for 74\% of all
  review-aspect mentions, and the top ten account for 96\%. A developer triaging in PageRank
  order completes three-quarters of the complaint volume after reading five aspect buckets,
  whereas random aspect ordering would require thousands of buckets to reach the same coverage,
  a roughly $10^3\times$ concentration ratio. The KG additionally enables per-aspect drill-down,
   distinguishing \texttt{battery}$\to$Samsung-drain from \texttt{battery}$\to$fast-charge.

\subsection{Ablations}\label{ablations}

  An ablation removes one component and re-measures, so that a component's contribution is
  demonstrated rather than assumed. Seven were designed and six were run; the seventh is deferred
  for a reason given below. Each one below states what it removes, what
  the removal is supposed to test, what happened, and what we conclude, because three of the
  seven came back against the design and a summary line would hide that.
  Table~\ref{tab:ablation-summary} collects the verdicts.

\begin{table}[t]\centering\footnotesize
\caption{Ablation summary. ``Removing this hurts'' is the only outcome that justifies a
component. Three of the seven do not clear that bar, and the paper treats them as findings
rather than as inconveniences.}\label{tab:ablation-summary}
\begin{tabular}{@{}p{0.9cm}p{3.6cm}p{3.6cm}p{2.3cm}@{}}\toprule
id & what is removed & effect of removing it & verdict \\\midrule
A1 & the KG, feeding flat clusters to Stage~3 & Stage~3 quality unchanged & no effect \\
A1b & the KG, against flat clustering matched to the same 605 clusters & flat wins on all three intrinsic metrics & \textbf{against the KG} \\
A2 & hierarchical sub-clustering & purity 0.625 flat vs 0.390 KG & \textbf{against the KG} \\
A3 & taxonomy grounding, so the LLM writes free-form & strict fill $0.96 \to 0.69$; per-type structure to zero & \textbf{justifies routing} \\
A4 & HITL-2 review of specs & no change to benchmark scores & no effect at benchmark time \\
A5 & retrieval, keeping the IssueSpec & all automatic metrics move $<1\%$ & \textbf{against RAG} \\
A6 & the IssueSpec, keeping retrieval & quality $4.59 \to 2.24$ & \textbf{justifies the IR} \\
A7 & single-stream RLHF feedback & not run & deferred \\
\bottomrule\end{tabular}\end{table}

  \textbf{A1, remove the knowledge graph and feed flat clusters into Stage~3.} This asks whether
  Stage~3 reads the graph or merely the group. Stage~3 quality does not change, because the
  translator consumes a cluster centroid and its representative reviews, not the graph edges. The
  KG is therefore not a Stage~3 dependency, which narrows where its value could lie.

  \textbf{A1b, compare the KG against flat clustering matched to the same number of clusters.}
  This is the reviewer-prompted version of A1 and it exists because of a confound. The KG
  produces 605 sub-clusters and the flat baseline produces 194, and intrinsic metrics move with
  cluster count on their own, so any comparison between the two conflates structure with
  granularity. Re-running agglomerative clustering forced to exactly 605 clusters removes the
  confound. The flat version then wins on all three metrics, Davies-Bouldin 1.12 against 2.24
  (lower is better), Calinski-Harabasz 4.21 against 1.85 (higher is better), and silhouette
  $+0.148$ against $-0.234$ (higher is better, and a negative value means the average point sits
  closer to a neighbouring cluster than to its own). The honest conclusion is that the KG does
  not produce geometrically better clusters. What it produces is an ordering, through PageRank
  centrality over the review-aspect graph, and prioritization is the only cluster-layer claim we
  make.

  \textbf{A2, remove hierarchical sub-clustering.} Weighted Y/P/N purity, judged on five
  representative reviews per cluster, is 0.625 for the flat 194 and 0.390 for the KG's 605. Finer
  granularity costs coherence here rather than buying it. Taken with A1b this is the second
  result against the cluster layer, and we report both.

  \textbf{A3, remove taxonomy grounding.} The LLM writes a free-form issue description instead of
  filling a routed template. Strict template-fill falls from 0.96 to 0.69, and the per-type
  structural columns fall to exactly zero, because an unrouted generator never emits a
  reproduction-step list or an As-a/I-want/So-that triple at all when nothing asks it to. This is
  the ablation that isolates the paper's central mechanism. The gain is not the language model,
  it is the routing.

  \textbf{A4, remove the HITL-2 spec review.} Benchmark scores do not move, and the reason is
  procedural rather than substantive. HITL-2 is a production-time gate; the 100 specs used in the
  evaluation were scored before any human correction, so the headline numbers already describe
  pre-HITL output. The ablation therefore says nothing about whether HITL-2 helps in deployment,
  and we do not claim that it does.

  \textbf{A5, remove retrieval and keep the IssueSpec.} BLEU-1 moves by $-0.006$, ROUGE-L by
  $-0.008$, and BERTScore F1 by $-0.004$, all under one percent absolute. Two caveats belong with
  this result. First, A5 was scored on automatic metrics only, while A6 was scored by humans, so
  the pair is not measured on one instrument and the contrast between them is directional rather
  than a like-for-like comparison. A human-rated no-retrieval condition is the missing cell.
  Second, this sits beside a finding that looks contradictory at first. Retrieval without the IR
  scores \emph{below} a no-context baseline, yet removing retrieval from the full system changes
  almost nothing. The reading that fits both is that unstructured retrieval pulls a generator
  toward generic precedent when nothing else constrains it, and that a typed specification in the
  prompt already constrains it, so the harm has nothing left to act on. What follows is not that
  retrieval helps, but that it stops mattering once the IR is present.

  \textbf{A6, remove the IssueSpec and keep retrieval.} This ablation is superseded. Under the
  confounded design it showed quality falling from 4.59 to 2.24; under the corrected same-model
  design (Section~\ref{exp2-results}) removing the IssueSpec changes quality by $-0.03$, which is
  not distinguishable from zero. A5 and A6 together now say something different from what we
  first concluded: neither retrieval nor the typed intermediate produces a measurable change in
  the reply a user receives, and the Stage~3 to Stage~4 coupling we claimed is not supported.

  \textbf{A7, single-stream RLHF feedback.} Not run. It was designed to test whether quality and
  compliance feedback can share one channel, which only becomes a meaningful question once the
  constraint binds and the two signals actually compete. Section~\ref{exp3-results} reaches the
  first half of that condition and not the second, so A7 stays deferred rather than reported.

%% =====================================================================
\section{Discussion}\label{discussion}

\subsection{What Carries the Result, and What Does Not}\label{load-bearing}
A five-stage pipeline invites the question of whether every stage earns its place. Our own
ablations say it does not, and separating the two groups is more useful to a reader than
defending the whole.

\emph{Load-bearing.} Two components carry the measured gains. Type-routed templating is the
principal driver of Stage~3 quality: removing taxonomy grounding (A3) drops strict template-fill
from 0.96 to 0.69, and the structural per-field metrics collapse to zero because an unrouted
generator does not emit the reproduction-step list or the As-a/I-want/So-that triple at all. The
the typed IR is \emph{not} a demonstrated driver of Stage~4 quality: under a same-model
comparison removing it moves quality by $0.03$ points ($p = 0.38$, Section~\ref{exp2-results}),
and removing retrieval instead (A5) moves every automatic metric by less than one point
absolute. Retrieval without the IR is significantly worse than a no-context
baseline, not merely equal to it. The intermediate carries the work, and the retrieval layer wrapped around
it does not.

\emph{Not load-bearing, and reported as such.} The knowledge graph does not win on intrinsic
cluster geometry. Under count matching (A1b), flat agglomerative clustering beats the KG on all
three intrinsic metrics, and sub-cluster purity is lower at the KG's finer granularity (A2).
The KG's defensible contribution is centrality-based prioritization, that is, which issues a
developer should see first, not tighter clusters; we make no cluster-quality claim for it. We
should be straight about how far even that goes. PageRank runs on the bipartite review-aspect
graph, and nothing in our results shows it needs the hierarchical sub-clustering layer to do so.
A flat clustering with the same aspect graph would plausibly support the same ranking, and we did
not test that configuration. Stage~2 as built is therefore justified in its prioritization role
and unjustified in its sub-clustering role, and a reader designing a lighter system should read
it that way.
Stage~5 sits in between. As Section~\ref{exp3-results} reports, the compliance constraint does
bind once the policy is capable of violating it, and the dual update engages rather than
collapsing, but no quality-versus-compliance trade-off is demonstrated at the scale we can train
at. It is released as a testbed with a partial result.
\textsc{SpecCov} falls in the same group and for a sharper reason: tested against human
faithfulness judgements it shows no correlation at all (Section~\ref{speccov-validation}), so we
withdraw it as a measure rather than defend it.

\emph{What this means for the framework claim.} The contribution that survives ablation is the
typed IR and the type-routed templating that produces it, together with the verified-anchor
procedure that makes the IR producible from noisy labels at corpus scale. Stages~2 and~5 are
supporting infrastructure with honestly mixed evidence. A reader who wants the smallest system that reproduces the headline result needs Stage~1, Stage~3,
and a generator conditioned on the resulting specification. That is less than Stage~4 as we built
it: A5 shows the retrieval layer inside Stage~4 contributes under one percent on automatic
metrics, so the part of Stage~4 that earns its place is the conditioning, not the retrieval.

\subsection{Implications}
Three claims travel beyond app reviews. \emph{First}, LLM annotation noise on
  software-engineering text is structurally biased rather than
  independent-and-identically-distributed: the 25\% mislabeling rate is concentrated on minority
   classes the LLM is blind to, not spread uniformly across the corpus. This is a finding about
  LLM-as-annotator failure modes that bears on any application using LLM labels for
  structured-text annotation. \emph{Second}, a typed structural intermediate, not retrieval
  volume or agentic iteration, is what does the load-bearing work in spec-aware response
  generation. That claim is withdrawn in Section~\ref{exp2-results}; what remains is the $+0.12$ from adding
  agentic critique-revise iteration to the same vanilla-RAG backbone, indicating the structural
  intermediate carries the work, not the iteration loop. \emph{Third}, the count-controlled
  cluster-validation ablation (A1b) delivers a self-contained methodological contribution:
  whenever a new clustering design changes cluster count and structure simultaneously, intrinsic
   geometry metrics (DB, CH, silhouette) conflate the two effects. The minimum reporting
  standard our analysis suggests is to (a)~compare against a count-matched baseline on the same
  intrinsic metrics, and (b)~pair that with a domain-utility metric (PageRank coverage in our
  case). This pattern applies directly to clustering-design comparisons in IR, software
  engineering, and NLP.

\subsection{Threats to Validity}
\emph{Construct.} Template-fill measures coverage rather than content correctness, and Section~\ref{speccov-validation} shows that \textsc{SpecCov} does not fill that gap: it is a lexical-grounding proxy that does not track human faithfulness judgement~\cite{grusky2018newsroom, maynez2020faithfulness, kryscinski2020evaluating, ji2023survey}. The paper therefore has no validated content-correctness measure for Stage~3, and the five-dimension rubric carries that burden alone. NLI- and QA-based faithfulness measurement~\cite{laban2022summac, honovich2022true, durmus2020feqa} is the clear next step.
\emph{Internal.} The count-versus-structure confound is exposed and resolved by the A1b ablation, which reframes C3's contribution as centrality-based prioritization rather than intrinsic geometry. A second internal threat we did not remove is that the 100 clusters carrying the Stage~3 and Stage~4 results were produced by the KG pipeline, the same layer A1b and A2 find wanting. Since flat clustering scores higher on purity, the headline gains may be measured on a slightly worse partition than a simpler system would produce, which would understate rather than inflate them. Re-running Stages~3 and~4 over flat clusters would settle it and is the cheapest open experiment in the paper.
\emph{External.} RRGen contains 58 anonymized Android apps; V5 transfers zero-shot to the Maalej label space. Stage~2--4 cross-corpus transfer remains to be validated. The headline LLM is Claude Opus, with Llama-3.3-70B and Qwen2.5-\{3B, 1.5B\} as cross-family bias controls; GPT-4o and Gemini-1.5-Pro replication is queued (Section~\ref{future-work}).
\emph{Conclusion.} The three-rater agreement on the classification gold ($\kappa = 0.968$ on the Y/N verdict) sits well above what free-labeling studies of app reviews report, which invites the question of whether the raters were genuinely independent. Three properties of the design account for the gap. The raters verify a supplied label rather than choose one from seven, which collapses most of the decision space; a 20-review calibration round with adjudicated disagreements precedes the main task; and the written protocol fixes eight boundary rules that carry the cases annotators most often split on. Residual disagreement is rater-specific rather than random, one rater accounts for most of the \texttt{usability} reassignments, which is the pattern independent annotation produces. The Stage~4 evaluation is now on the same footing: all three raters scored the complete 400-row set, and the headline gain reproduces individually for each of them.
\emph{Reliability, and which $\alpha$ belongs to what.} Three agreement statistics appear in this paper and they measure different things, so a single ``moderate $\alpha$'' reading would be wrong. The classification gold is annotated by three humans and reaches Fleiss $\kappa = 0.968$ on the verdict and $0.979$ on the seven-class label. The Stage-4 response set is rated by the same three humans on all 400 rows, with quality weighted $\kappa$ from 0.970 to 0.986 and Krippendorff $\alpha = 0.781$ on the harder binary helpful judgement. The one moderate figure, $\alpha = 0.451$, comes from a different and deliberately weaker panel: a cluster-purity check whose three raters are the lead author plus two Qwen2.5-3B prompts, following the LLM-as-rater convention of Gilardi et al.~\cite{gilardi2023chatgpt}. It bounds how far an LLM rater tracks a human on a subjective sub-theme judgement; it is not the reliability of any headline result, and no claim in this paper rests on it alone.
\emph{Statistical.} Experiment~2's four simultaneous tests reject under non-overlapping nulls and survive Bonferroni correction at $\alpha = 0.05$.

\subsection{Reproducibility and Extensibility}\label{reproducibility}
Every numerical claim in Section~\ref{results} is backed by a saved data file and runs end-to-end from raw RRGen to Stage~4 output in approximately six hours on a single GPU-equipped workstation. The release bundles:
\begin{itemize}\setlength{\itemsep}{1pt}
\item V1--V5 classifier checkpoints + cleanlab correction procedure;
\item the 5{,}230-review verified anchor and 490-review classification gold;
\item the 400-row blinded evaluation and the count-controlled A1b output;
\item the \textsc{SpecCov} scorer plus the 60-spec, three-rater human faithfulness ratings and the validation analysis;
\item the CMDP-RLHF testbed (KTO, DPO, constrained-proxy, and Lagrangian Constrained PPO trainers);
\item the 100-cluster benchmark, with the 99-review LLM-assisted purity subsample ($\alpha = 0.451$) kept separate from the all-human agreement sets;
\item a pinned-environment Docker image (single-command build on any CUDA~12.x host);
\item an eighteen-segment verifier that recomputes every numerical claim in Section~\ref{results}
from the saved data files, including the three-rater agreement statistics, the cross-family
rubric, the \textsc{SpecCov} validation, the binding-constraint run, and the generation-variance
sweep, so a reader can check the paper against the artifact without re-running the pipeline.
\end{itemize}
The pipeline is composable by construction: because each stage emits a typed intermediate, components (classifier, clusterer, taxonomy template, alignment trainer) can be substituted independently without breaking downstream contracts. Future researchers can adopt or replace any single component while keeping the rest of the IR-centered evaluation harness intact. Deterministic components (clustering, classifier training) use fixed seeds; the LLM-dependent stages (annotation, IR generation, and response generation) reproduce statistically rather than bit-identically. We measured how large that sampling variance is. Regenerating Stage~3 specs three times over the same 15 stratified clusters with the local Qwen2.5-3B generator at temperature 0.7, changing only the seed, gives strict template-fill 0.685, 0.700, and 0.634 (mean 0.673, sd 0.029, range 0.067) and \textsc{SpecCov} 2.73, 2.67, and 3.00 (mean 2.80, sd 0.144, range 0.333), with no JSON parse failures in any run. Run-to-run noise on template-fill is therefore around three points on a 0--1 scale for this generator and prompt. That is the yardstick a reader should hold up against the differences we report. The gaps the paper rests on clear it comfortably: 0.96 versus 0.69 for taxonomy grounding, 0.96 versus 0.48 across the model families, and $+2.35$ Likert points for the IssueSpec. Differences of a few hundredths between adjacent conditions, such as the 4.09 versus 3.96 at the top of the cross-family rubric, do not clear it and should not be read as a stable ordering on their own. Two limits remain: the headline Claude generations were not repeated, because that needs API budget rather than new code, and Stage~4 response generation was not repeated either. The released script runs both sweeps unchanged.

\subsection{Future Work}\label{future-work}
Five extensions map directly to the limitations above; each starts from a resource released with this paper.

\textbf{(0) Replace \textsc{SpecCov} and widen the cross-family comparison.} The
\textsc{SpecCov} validation in Section~\ref{speccov-validation} came back negative, so Stage~3
needs a faithfulness measure that actually works. The natural replacement is entailment-based:
score each spec field against the source reviews with an NLI model~\cite{laban2022summac,
honovich2022true} or a question-generation-and-answering pipeline~\cite{durmus2020feqa}, then
validate that measure against the same 60-spec human ratings we release here, which now serve as
a ready-made benchmark. Second, the cross-family rubric comparison
(Table~\ref{tab:xfam-rubric}) rests on 14 matched bug-report clusters for the four-way panel,
because that is the extent of the overlap between the four generators' spec sets, and it is
scored by a judge from one of the four families. Generating the missing spec types for the two
Qwen models and re-scoring with a judge outside all four families would remove both limits.
Both are runs rather than redesigns.

\textbf{(1) Scale the benchmark.} The Stage~4 set now carries three annotators on all 400 rows; the remaining step is to grow it from 100 clusters to 200--300 and to put three annotators on each IssueSpec as well, reporting Krippendorff's $\alpha$ at the IssueSpec level.

\textbf{(2) CMDP RLHF at a scale where the trade-off shows.} The binding-constraint re-run in
Section~\ref{exp3-results} establishes that the dual objective engages; what it does not show is
a trade-off, because a 5\% base violation rate over 30 steps supplies too little signal. The next
run wants a longer horizon and a policy with a higher base violation rate, on an 8B-class model.
\emph{Original plan.} Where the compliance constraint is expected to bind, activating the Lagrange multiplier and making the Pareto-dominance comparison meaningful.

\textbf{(3) Frontier-family replication.} GPT-4o and Gemini-1.5-Pro on Stages~3--4. The current paper closes the Meta-family slice with Llama-3.3-70B and the Alibaba slice with Qwen2.5-\{3B, 1.5B\}.

\textbf{(4) Cross-corpus transfer.} Stages~2--4 on additional review corpora; V5 already transfers zero-shot to Maalej~\cite{maalej2015bug}.

\textbf{(5) Agentic Planner extension to Stage~4.} The Planner dispatches a different
  sub-agent profile per issue type: defect-repair for \texttt{bug\_report} via
  SWE-Agent~\cite{yang2024swe} or RepairAgent~\cite{bouzenia2025repairagent}, scaffolding
  for \texttt{feature\_request}, profiler-driven optimization for \texttt{performance}, a
  Nielsen-heuristic checker for \texttt{usability}, and a device-OS test-matrix executor for
  \texttt{compatibility}. A spec-level stub already runs the dispatch across all five issue
  types; end-to-end patch evaluation on real repositories through
  HyperAgent~\cite{phan2024hyperagent}, measured by patch acceptance and test-pass rate, is the
  next step.

%% =====================================================================
\section{Conclusion}\label{conclusion}

 
  We proposed \textbf{IssueSpec}, a knowledge-graph-grounded typed intermediate representation
  that closes the five existing research gaps identified in Section~\ref{introduction}:
  classification stops at the label, response generation skips the structured intermediate,
  agentic software engineering assumes the spec exists, knowledge graphs emit no typed schema,
  and RLHF has not addressed dev-rel compliance. An isolation experiment shows the typed IR is
  the load-bearing component of spec-aware response generation: adding IssueSpec to RAG raises
  human-rated quality by $+0.03$ Likert points, pooled over three independent raters on the corrected same-model comparison ($p = 0.38$), while
  retrieval without the IR scores significantly below a no-context baseline. A
  verified-anchor procedure makes the IR producible at corpus scale from noisy LLM labels,
  lifting Cohen's $\kappa$ from 0.163 to 0.592, and to 0.616 on the strictly held-out gold, with 88.66\% classifier endorsement
  of corrections. Cross-family replication on four LLMs (Claude, Llama, Qwen-3B, Qwen-1.5B)
  preserves the rank order, ruling out single-model bias. Two findings travel beyond app
  reviews: LLM annotation noise on software-engineering text is structurally biased rather than
  random, and a typed structural intermediate is what does the load-bearing work in spec-aware
  response generation. We release the IssueSpec schema, the verified anchor, the
  \textsc{SpecCov} scorer together with the human ratings that show it does not measure
  faithfulness, the count-controlled cluster-validation ablation, the
  structured-RAG architecture, the CMDP-RLHF trainers, and the benchmark dataset to support
  follow-on research.

%% =====================================================================
\section*{GenAI Usage Disclosure}
We disclose all uses of generative AI in this work in line with Elsevier policy on the use of generative AI in scientific writing. \emph{Writing and code assistance.} Anthropic's Claude (Opus 4.7)~\cite{anthropic2025claude} assisted with grammatical refinement, caption polishing, LaTeX formatting, and reviewing implementation patterns and statistical-test code. All pipeline code, experimental scripts, and analysis were written by the human authors. Every structural, framing, experimental-design, and claim decision was made by the authors, and every reported number was produced by author-written code run under human supervision and verified against saved data files. \emph{Models inside the system.} Several generative models are themselves pipeline components and are reported as such in the body: Claude (Opus 4.7) is the headline Stage~3 generator (Table~\ref{tab:stage3}); Llama-3.3-70B-Instruct (via Groq) and Qwen2.5-\{3B, 1.5B\}-Instruct are cross-family Stage~3 comparators (Table~\ref{tab:multi-llm}); Qwen2.5-3B-Instruct is the LLM-as-judge for the five-dimension rubric and the Y/P/N cluster-purity audit, both calibrated against the lead author (Tables~\ref{tab:rubric}, \ref{tab:cluster-quality}). The V2 corpus labels in Section~\ref{exp1-results} are LLM-generated and explicitly treated as
   noisy, which is precisely what motivates the verified-anchor procedure of Stage~1. \emph{Human verification.} No reported numerical result, expert annotation, or ground-truth label was automatically generated by AI without human verification. The 5{,}230-review verified anchor, the 490-review classification gold, the 99-review three-rater purity ratings, and the 20 lead-author IssueSpec references were all produced and reviewed by human annotators (Section~\ref{experiments}).

%% =====================================================================
\section*{CRediT authorship contribution statement}
\textbf{Fabiha Jalal:} Conceptualization, Methodology, Software, Validation, Formal analysis,
Investigation, Data curation, Writing -- original draft, Writing -- review and editing,
Visualization. \textbf{Sadia Tasnim Dhruba:} Investigation, Data curation, Validation, Writing
-- review and editing. \textbf{Tazrian Zaman Tushti:} Investigation, Data curation, Validation,
Writing -- review and editing. \textbf{Ajwad Abrar:} Methodology, Supervision, Writing --
review and editing. \textbf{Md Kamrul Hasan:} Supervision, Resources, Writing -- review and
editing. \textbf{Hasan Mahmud:} Supervision, Resources, Project administration, Writing --
review and editing.

\section*{Declaration of competing interest}
The authors declare that they have no known competing financial interests or personal
relationships that could have appeared to influence the work reported in this paper.

\section*{Data availability}
All code, configuration files, and analysis scripts are available in the project repository.
A verification bundle that reproduces every numerical claim in this paper is archived on Zenodo
(DOI 10.5281/zenodo.21982774, v2.0), and the Stage-1 V5 classifier is released on Hugging Face.
That bundle supersedes an earlier release (DOI 10.5281/zenodo.20320410) which predates the
Stage-4 correction described in Section~\ref{exp2-results}; both versions remain resolvable, and
the earlier one should not be used.
The underlying review corpus derives from the publicly released RRGen dataset of Gao et al.;
derived artifacts are redistributed subject to the terms of that original release.

\section*{Acknowledgements}
We thank the human annotators who produced the verified anchor and the lead-author reference
set. The authors verified all technical content, experimental results, and claims in this paper.

%% =====================================================================
\bibliographystyle{elsarticle-num}
\bibliography{references}

\appendix
  
  \section{Implementation Settings}\label{appendix:settings}

  Section~\ref{system-design} states what each stage does and why. This appendix collects the
  exact settings, so the method section carries the logic and the numbers stay checkable.
  Stage~1 fine-tunes a RoBERTa multi-label head; confident-learning correction against the
  verified anchor is gated by anchor confidence $\geq 0.70$ and LLM probability $\leq 0.20$,
  with a $3\times3$ sensitivity grid over $(0.65, 0.70, 0.75) \times (0.15, 0.20, 0.25)$
  reported in Section~\ref{exp1-results}. The \texttt{compatibility} stratum is seeded with 200
  template-generated synthetic reviews (device, OS, and screen templates) plus 100 keyword-mined
  RRGen reviews, all human-checked. Stage~2 embeds with Sentence-BERT, reduces with UMAP, and
  clusters with HDBSCAN, aspect-first, producing 605 sub-clusters of mean size 3.0 against 194
  flat clusters of mean size 51.5; PageRank runs on the bipartite review-aspect graph. Stage~3
  routes to one of five templates and re-prompts once with dimension-level feedback on any
  rubric failure. Stage~4 retrieves over a 15{,}100-document ChromaDB index built from five
  sources, with the IssueSpec injected per query rather than pre-indexed. Stage~5 trains five
  policies on distilGPT2 (82M parameters) from 400 SFT samples, 296 KTO pairs, and 100 DPO
  pairs; the Lagrangian loop runs 30 steps at batch size 4 with KL coefficient 0.05, and the
  compliance threshold sweep covers $\tau \in \{0.0, 0.5, 0.7, 0.9, 1.0\}$.
  \textsc{SpecCov} counts substantive-token overlap with the source cluster, tokens of length
  $\geq 5$ after stop-word filtering, with quoted-string and coverage bonuses.

\section{Annotator Guidelines}\label{appendix:annotators}

  The lead author (PhD in software engineering with prior mobile-app review-mining experience)
  served as primary annotator across all five annotation tasks. The 490-review classification
  gold was additionally annotated by two further independent human raters on the identical item
  set, aligned by review id; all three first completed a shared 20-review calibration round,
  disagreements were adjudicated before the main task opened, and the released gold label is the
  majority vote over the three raters. All human annotation in this work is carried out by
  those same three people. One of the two additional raters provided inter-rater validation on a
  30-pair Stage~4 subsample, and two Qwen2.5-3B-Instruct prompts (concise role-based and
  chain-of-thought) served as the third and fourth raters in the 99-review Krippendorff $\alpha$
   calibration. For the 5{,}230-review verified anchor, annotators classified each review into
  one of seven Maalej-Nabil classes (\texttt{bug\_report}, \texttt{feature\_request},
  \texttt{performance}, \texttt{usability}, \texttt{compatibility}, \texttt{praise},
  \texttt{other}) with multi-label assignment permitted when a review carried multiple concerns;
   reviews below the ``fairly confident'' threshold were skipped and excluded from the anchor.
  For the 490-review classification gold standard, a stratified sample of 70 reviews per class
  was drawn from the working corpus, and the lead author re-labeled each review independently
  and blinded to the V2 LLM prediction; the 490 gold labels serve as the held-out evaluation set
   for Cohen's $\kappa$ progression (Table~\ref{tab:kappa}) and per-class F1
  (Table~\ref{tab:perclass-f1}). For the 5-dimension IssueSpec rubric (all 100 specs; an earlier stratified $n=28$ subsample, 6 per issue
   type plus 2 extras), each spec was scored on a 1--5 Likert scale across five dimensions:
  \emph{completeness} (are all required template fields filled?), \emph{accuracy} (do field
  values correctly reflect the source cluster?), \emph{actionability} (would a developer know
  what to do?), \emph{specificity} (does the spec name concrete entities rather than generic
  terms?), and \emph{clarity} (is the spec easy to read and unambiguous?), with each dimension
  scored independently. For the 400-row Stage~4 response evaluation (100 reviews $\times$ 4
  conditions: rrgen, prompt-baseline, \texttt{no\_spec}, \texttt{full}), the lead author rated
  \emph{quality} (1--5), \emph{specificity} (1--5), and \emph{helpful} (Y/N) while blinded to
  the generation method; condition labels were replaced with random A/B/C/D codes and unsealed
  only at aggregation, with the assignment persisted as spreadsheets with JSON master keys for
  audit. A second human rater independently scored a 30-pair (\texttt{no\_spec}, \texttt{full})
  subset for inter-rater calibration. For the Y/P/N cluster-purity rubric (initial 50-cluster
  audit; 99-cluster three-rater subsample), annotators read 5 representative reviews per cluster
   and assigned \emph{Y} (all 5 share the same sub-theme within the cluster's aspect), \emph{P}
  (3 or 4 of 5 share the sub-theme), or \emph{N} (fewer than 3 share); weighted purity is
  $\mathrm{purity}_w = (|Y| + 0.5|P|)/(|Y|+|P|+|N|)$, with $Y$ contributing 1.0, $P$
  contributing 0.5, and $N$ contributing 0. All comparative evaluations were blinded throughout
  the rating window and unsealed only at aggregation; reviews or specs below the annotator
  confidence threshold were marked as ``skip'' and excluded from downstream analysis.


\end{document}
